
\documentclass{article}
\usepackage{colortbl}
\usepackage{makecell}
\usepackage{multirow}
\usepackage{supertabular}

\begin{document}

\newcounter{utterance}

\centering \large Interaction Transcript for game `clean\_up', experiment `0\_easy\_7obj\_ga', episode 0 with Qwen3.6{-}35B{-}A3B{-}FP8{-}without{-}reasoning.
\vspace{24pt}

{ \footnotesize  \setcounter{utterance}{1}
\setlength{\tabcolsep}{0pt}
\begin{supertabular}{c@{$\;$}|p{.15\linewidth}@{}p{.15\linewidth}p{.15\linewidth}p{.15\linewidth}p{.15\linewidth}p{.15\linewidth}}
    \# & $\;$A & \multicolumn{4}{c}{Game Master} & $\;\:$B\\
    \hline

    \theutterance \stepcounter{utterance}  
    & & \multicolumn{4}{p{0.6\linewidth}}{
        \cellcolor[rgb]{0.9,0.9,0.9}{
            \makecell[{{p{\linewidth}}}]{
                \texttt{\tiny{[P1$\langle$GM]}}
                \texttt{Is mise do mháistir cluiche, agus tá tú ag imirt cluiche comhoibríoch leis an ngreille seo a leanas mar chlár cluiche:} \\
\\ 
\texttt{```} \\
\texttt{ 1234567} \\
\texttt{╔═╤═══╤═╗} \\
\texttt{║◌│◌◌◌│◌║ 1} \\
\texttt{╟─┼──┬┘◌║ 2} \\
\texttt{║◌│◌◌│S◌║ 3} \\
\texttt{╟─┴──┘DT║ 4} \\
\texttt{║◌O◌◌◌◌◌║ 5} \\
\texttt{║◌◌◌◌◌◌◌║ 6} \\
\texttt{║◌◌APM◌◌║ 7} \\
\texttt{╚═══════╝} \\
\\ 
\texttt{```} \\
\\ 
\texttt{* Taispeánann an imeall uachtarach comhordanáidí x ag méadú i dtreo na láimhe deise, agus an t{-}imeall ar dheis comhordanáidí y ag méadú síos.} \\
\texttt{* Cuirtear na rudaí seo a leanas go randamach ar do ghreille: 'P', 'O', 'T', 'S', 'D', 'A', 'M'.} \\
\\ 
\\ 
\texttt{Feiceann an t{-}imreoir eile leagan éagsúil den chlár cluiche, áit a gcuirtear na rudaí in áiteanna randamacha éagsúla. Ní féidir leat clár an imreora eile a fheiceáil, agus ní féidir leo do cheannsa a fheiceáil.} \\
\\ 
\\ 
\texttt{**Sprioc:**} \\
\\ 
\texttt{Caithfidh an dá imreoir na rudaí a bhogadh ar a gcúlra féin ionas go gcríochnóidh rudaí comhionanna ag na comhordanáidí céanna. Caithfidh tú cumarsáid a dhéanamh leis an imreoir eile chun cumraíocht sprice choiteann a aontú.} \\
\\ 
\\ 
\texttt{**Rialacha:**} \\
\\ 
\texttt{* I ngach babhta, is féidir leat ceann amháin go díreach den dá ordú seo a leanas a sheoladh:} \\
\texttt{ 1. `ABAIR: <TEACHTAIREACHT>`: chun teachtaireacht (gach rud suas go dtí an chéad bhriseadh líne eile) a sheoladh chuig an imreoir eile. Cuirfidh mé ar aghaidh í chuig do pháirtí.} \\
\texttt{2. `BOG: <RUD>, (<X>, <Y>)`: chun rud a bhogadh go suíomh nua, áit a bhfuil `<X>` mar an colún agus `<Y>` mar an ró. Cuirfidh mé in iúl duit an raibh do bhogadh rathúil nó nach raibh.} \\
\texttt{* Mura gcloíonn tú leis an bhformáid, nó má sheolann tú roinnt orduithe ag an am céanna, caithfidh mé pionós a ghearradh ort.} \\
\texttt{* Má charnaíonn an dá imreoir níos mó ná 16 pionóis, caillfidh sibh beirt an cluiche.} \\
\texttt{* Tá sé ríthábhachtach go ndéanann tú cumarsáid leis an imreoir eile maidir le do staid sprice! Is é an t{-}aon bhealach *amháin* ar féidir leat do straitéis a chur ar aghaidh chuig an imreoir eile ná an t{-}ordú `ABAIR: <TEACHTAIREACHT>` a úsáid!} \\
\\ 
\\ 
\texttt{**Rudaí a Bhogadh**} \\
\\ 
\texttt{* Ní féidir leat rudaí a bhogadh ach chuig cealla laistigh de theorainneacha na greille. Caithfidh an chill sprice a bheith folamh, is é sin, ní mór di ach an tsiombail '◌' a bheith inti.} \\
\texttt{* Má dhéanann tú iarracht rud a bhogadh go spota nach bhfuil folamh, nó má dhéanann tú iarracht é a bhogadh lasmuigh den ghreille, caithfidh mé pionós a ghearradh ort. Gheobhaidh tú triail eile.} \\
\texttt{* Sula ndéanann tú bogadh, seiceáil faoi dhó go bhfuil an spota sprice folamh, agus nach bhfuil aon litir, fráma ná líne ann!} \\
\\ 
\\ 
\texttt{**Deireadh an Chluiche**} \\
\\ 
\texttt{Má cheapann tú gur shroich tú an sprioc maidir le gach rud a ailíniú, is féidir leat iarraidh ar an imreoir eile an cluiche a chríochnú trí `ABAIR: críochnaithe?` a sheoladh. Má iarrann an t{-}imreoir eile ort an cluiche a chríochnú, agus má fhreagraíonn tú le `ABAIR: críochnaithe!`, tiocfaidh deireadh leis an gcluiche.} \\
\\ 
\texttt{Buaíonn an dá imreoir má chríochnaíonn an cluiche laistigh de 28 babhta, áit a shainmhínítear babhta amháin mar bheirt imreoirí ag seoladh ordú bailí an duine.} \\
\\ 
\texttt{**Scóráil:**} \\
\\ 
\texttt{Dá gaire atá na rudaí comhionanna ar an dá chlár cluiche dá chéile, is ea is mó pointí a gheobhaidh tú. Laghdaíonn pionóis do phointí. An féidir leat an taifead a shárú?} \\
\\ 
\\ 
\texttt{Seol teachtaireacht chuig an imreoir eile chun an cluiche a thosú!} \\
            }
        }
    }
    & \\ \\

    \theutterance \stepcounter{utterance}  
    & \multicolumn{4}{p{0.6\linewidth}}{
        \cellcolor[rgb]{0.9,0.9,0.9}{
            \makecell[{{p{\linewidth}}}]{
                \texttt{\tiny{[P1$\rangle$GM]}}
                \texttt{ABAIR: Tosaímid an cluiche! Seo mo staid reatha:} \\
\texttt{P ag (1,7)} \\
\texttt{O ag (2,5)} \\
\texttt{T ag (4,4)} \\
\texttt{S ag (6,3)} \\
\texttt{D ag (5,4)} \\
\texttt{A ag (6,7)} \\
\texttt{M ag (7,7)} \\
\\ 
\texttt{Cad iad do chomhordanáidí do rudaí féin? Seol iad chugam chun go mbeidh muid comhoiriúnach.} \\
            }
        }
    }
    & & \\ \\

    \theutterance \stepcounter{utterance}  
    & & & \multicolumn{2}{p{0.3\linewidth}}{
        \cellcolor[rgb]{0.9,0.9,0.9}{
            \makecell[{{p{\linewidth}}}]{
                \texttt{\tiny{[GM$|$GM]}}
                \texttt{ABAIR: Tosaímid an cluiche! Seo mo staid reatha:} \\
\texttt{P ag (1,7)} \\
\texttt{O ag (2,5)} \\
\texttt{T ag (4,4)} \\
\texttt{S ag (6,3)} \\
\texttt{D ag (5,4)} \\
\texttt{A ag (6,7)} \\
\texttt{M ag (7,7)} \\
\\ 
\texttt{Cad iad do chomhordanáidí do rudaí féin? Seol iad chugam chun go mbeidh muid comhoiriúnach.} \\
            }
        }
    }
    & & \\ \\

    \theutterance \stepcounter{utterance}  
    & & & \multicolumn{2}{p{0.3\linewidth}}{
        \cellcolor[rgb]{0.9,0.9,0.9}{
            \makecell[{{p{\linewidth}}}]{
                \texttt{\tiny{[GM$|$GM]}}
                \texttt{Invalid format: tail is not empty: '} \\
\texttt{P ag (1,7)} \\
\texttt{O ag (2,5)} \\
\texttt{T ag (4,4)} \\
\texttt{S ag (6,3)} \\
\texttt{D ag (5,4)} \\
\texttt{A ag (6,7)} \\
\texttt{M ag (7,7)} \\
\\ 
\texttt{Cad iad do chomhordanáidí do rudaí féin? Seol iad chugam chun go mbeidh muid comhoiriúnach.} \\
            }
        }
    }
    & & \\ \\

    \theutterance \stepcounter{utterance}  
    & & \multicolumn{4}{p{0.6\linewidth}}{
        \cellcolor[rgb]{0.9,0.9,0.9}{
            \makecell[{{p{\linewidth}}}]{
                \texttt{\tiny{[P1$\langle$GM]}}
                \texttt{Pionós: Ní cheadaítear aon rud a bheith i do theachtaireacht tar éis an ordaithe!} \\
\\ 
\texttt{Cinntigh nach bhfuil i do fhreagra ach `ABAIR: <TEACHTAIREACHT>` nó `BOG: <RUD>, (<X>, <Y>)`, agus tada eile!} \\
\\ 
\texttt{Tá 1 as 16 pionós carntha agaibh le chéile.} \\
\texttt{Déan iarracht arís, le do thoil!} \\
            }
        }
    }
    & \\ \\

    \theutterance \stepcounter{utterance}  
    & \multicolumn{4}{p{0.6\linewidth}}{
        \cellcolor[rgb]{0.9,0.9,0.9}{
            \makecell[{{p{\linewidth}}}]{
                \texttt{\tiny{[P1$\rangle$GM]}}
                \texttt{ABAIR: Mo rudaí: P(1,7), O(2,5), T(4,4), S(6,3), D(5,4), A(6,7), M(7,7). Cad iad do chomhordanáidí?} \\
            }
        }
    }
    & & \\ \\

    \theutterance \stepcounter{utterance}  
    & & & \multicolumn{2}{p{0.3\linewidth}}{
        \cellcolor[rgb]{0.9,0.9,0.9}{
            \makecell[{{p{\linewidth}}}]{
                \texttt{\tiny{[GM$|$GM]}}
                \texttt{ABAIR: Mo rudaí: P(1,7), O(2,5), T(4,4), S(6,3), D(5,4), A(6,7), M(7,7). Cad iad do chomhordanáidí?} \\
            }
        }
    }
    & & \\ \\

    \theutterance \stepcounter{utterance}  
    & & \multicolumn{4}{p{0.6\linewidth}}{
        \cellcolor[rgb]{0.9,0.9,0.9}{
            \makecell[{{p{\linewidth}}}]{
                \texttt{\tiny{[GM$\rangle$P2]}}
                \texttt{Is mise do mháistir cluiche, agus tá tú ag imirt cluiche comhoibríoch leis an ngreille seo a leanas mar chlár cluiche:} \\
\\ 
\texttt{```} \\
\texttt{ 1234567} \\
\texttt{╔═╤═══╤═╗} \\
\texttt{║◌│◌◌◌│◌║ 1} \\
\texttt{╟─┼──┬┘◌║ 2} \\
\texttt{║◌│◌M│◌T║ 3} \\
\texttt{╟─┴──┘P◌║ 4} \\
\texttt{║D◌◌◌◌A◌║ 5} \\
\texttt{║◌◌◌◌O◌◌║ 6} \\
\texttt{║◌◌S◌◌◌◌║ 7} \\
\texttt{╚═══════╝} \\
\\ 
\texttt{```} \\
\\ 
\texttt{* Taispeánann an imeall uachtarach comhordanáidí x ag méadú i dtreo na láimhe deise, agus an t{-}imeall ar dheis comhordanáidí y ag méadú síos.} \\
\texttt{* Cuirtear na rudaí seo a leanas go randamach ar do ghreille: 'P', 'O', 'T', 'S', 'D', 'A', 'M'.} \\
\\ 
\\ 
\texttt{Feiceann an t{-}imreoir eile leagan éagsúil den chlár cluiche, áit a gcuirtear na rudaí in áiteanna randamacha éagsúla. Ní féidir leat clár an imreora eile a fheiceáil, agus ní féidir leo do cheannsa a fheiceáil.} \\
\\ 
\\ 
\texttt{**Sprioc:**} \\
\\ 
\texttt{Caithfidh an dá imreoir na rudaí a bhogadh ar a gcúlra féin ionas go gcríochnóidh rudaí comhionanna ag na comhordanáidí céanna. Caithfidh tú cumarsáid a dhéanamh leis an imreoir eile chun cumraíocht sprice choiteann a aontú.} \\
\\ 
\\ 
\texttt{**Rialacha:**} \\
\\ 
\texttt{* I ngach babhta, is féidir leat ceann amháin go díreach den dá ordú seo a leanas a sheoladh:} \\
\texttt{ 1. `ABAIR: <TEACHTAIREACHT>`: chun teachtaireacht (gach rud suas go dtí an chéad bhriseadh líne eile) a sheoladh chuig an imreoir eile. Cuirfidh mé ar aghaidh í chuig do pháirtí.} \\
\texttt{2. `BOG: <RUD>, (<X>, <Y>)`: chun rud a bhogadh go suíomh nua, áit a bhfuil `<X>` mar an colún agus `<Y>` mar an ró. Cuirfidh mé in iúl duit an raibh do bhogadh rathúil nó nach raibh.} \\
\texttt{* Mura gcloíonn tú leis an bhformáid, nó má sheolann tú roinnt orduithe ag an am céanna, caithfidh mé pionós a ghearradh ort.} \\
\texttt{* Má charnaíonn an dá imreoir níos mó ná 16 pionóis, caillfidh sibh beirt an cluiche.} \\
\texttt{* Tá sé ríthábhachtach go ndéanann tú cumarsáid leis an imreoir eile maidir le do staid sprice! Is é an t{-}aon bhealach *amháin* ar féidir leat do straitéis a chur ar aghaidh chuig an imreoir eile ná an t{-}ordú `ABAIR: <TEACHTAIREACHT>` a úsáid!} \\
\\ 
\\ 
\texttt{**Rudaí a Bhogadh**} \\
\\ 
\texttt{* Ní féidir leat rudaí a bhogadh ach chuig cealla laistigh de theorainneacha na greille. Caithfidh an chill sprice a bheith folamh, is é sin, ní mór di ach an tsiombail '◌' a bheith inti.} \\
\texttt{* Má dhéanann tú iarracht rud a bhogadh go spota nach bhfuil folamh, nó má dhéanann tú iarracht é a bhogadh lasmuigh den ghreille, caithfidh mé pionós a ghearradh ort. Gheobhaidh tú triail eile.} \\
\texttt{* Sula ndéanann tú bogadh, seiceáil faoi dhó go bhfuil an spota sprice folamh, agus nach bhfuil aon litir, fráma ná líne ann!} \\
\\ 
\\ 
\texttt{**Deireadh an Chluiche**} \\
\\ 
\texttt{Má cheapann tú gur shroich tú an sprioc maidir le gach rud a ailíniú, is féidir leat iarraidh ar an imreoir eile an cluiche a chríochnú trí `ABAIR: críochnaithe?` a sheoladh. Má iarrann an t{-}imreoir eile ort an cluiche a chríochnú, agus má fhreagraíonn tú le `ABAIR: críochnaithe!`, tiocfaidh deireadh leis an gcluiche.} \\
\\ 
\texttt{Buaíonn an dá imreoir má chríochnaíonn an cluiche laistigh de 28 babhta, áit a shainmhínítear babhta amháin mar bheirt imreoirí ag seoladh ordú bailí an duine.} \\
\\ 
\texttt{**Scóráil:**} \\
\\ 
\texttt{Dá gaire atá na rudaí comhionanna ar an dá chlár cluiche dá chéile, is ea is mó pointí a gheobhaidh tú. Laghdaíonn pionóis do phointí. An féidir leat an taifead a shárú?} \\
\\ 
\texttt{Thosaigh an t{-}imreoir eile an cluiche trí an teachtaireacht seo a sheoladh:} \\
\\ 
\texttt{"Mo rudaí: P(1,7), O(2,5), T(4,4), S(6,3), D(5,4), A(6,7), M(7,7). Cad iad do chomhordanáidí?"} \\
\\ 
\texttt{Cad é do chéad ordú?} \\
            }
        }
    }
    & \\ \\

    \theutterance \stepcounter{utterance}  
    & & & \multicolumn{4}{p{0.6\linewidth}}{
        \cellcolor[rgb]{0.9,0.9,0.9}{
            \makecell[{{p{\linewidth}}}]{
                \texttt{\tiny{[GM$\langle$P2]}}
                \texttt{ABAIR: Mo rudaí: P(4,4), O(5,6), T(6,3), S(3,7), D(1,5), A(6,5), M(4,3).} \\
            }
        }
    }
     \\ \\

    \theutterance \stepcounter{utterance}  
    & & & \multicolumn{2}{p{0.3\linewidth}}{
        \cellcolor[rgb]{0.9,0.9,0.9}{
            \makecell[{{p{\linewidth}}}]{
                \texttt{\tiny{[GM$|$GM]}}
                \texttt{ABAIR: Mo rudaí: P(4,4), O(5,6), T(6,3), S(3,7), D(1,5), A(6,5), M(4,3).} \\
            }
        }
    }
    & & \\ \\

    \theutterance \stepcounter{utterance}  
    & & \multicolumn{4}{p{0.6\linewidth}}{
        \cellcolor[rgb]{0.9,0.9,0.9}{
            \makecell[{{p{\linewidth}}}]{
                \texttt{\tiny{[P1$\langle$GM]}}
                \texttt{Cuireadh do theachtaireacht ar aghaidh chuig an imreoir eile.} \\
\\ 
\\ 
\texttt{Tá sibh ag imirt babhta 1 faoi láthair as uasmhéid 28.} \\
\texttt{Tá 1 as 16 pionós carntha agaibh le chéile.} \\
\texttt{Sheol an t{-}imreoir eile an teachtaireacht seo:} \\
\\ 
\\ 
\texttt{"Mo rudaí: P(4,4), O(5,6), T(6,3), S(3,7), D(1,5), A(6,5), M(4,3)."} \\
\\ 
\\ 
\texttt{Cad é do chéad ordú eile?} \\
            }
        }
    }
    & \\ \\

    \theutterance \stepcounter{utterance}  
    & \multicolumn{4}{p{0.6\linewidth}}{
        \cellcolor[rgb]{0.9,0.9,0.9}{
            \makecell[{{p{\linewidth}}}]{
                \texttt{\tiny{[P1$\rangle$GM]}}
                \texttt{ABAIR: Go maith! Anois tá a fhios againn gach suíomh. Chun an cluiche a chríochnú go héifeachtach, ní mór dúinn na rudaí céanna a bhogadh go dtí na comhordanáidí céanna ar an dá thaobh.} \\
\\ 
\texttt{Seo an plean atá agam:} \\
\texttt{1. Ailínímid **P** ag (1,7) ar an dá thaobh. (Tá sé agam ag (1,7), tá sé aige ag (4,4). Caithfimid é a bhogadh go dtí suíomh comhionann. Ach fan, ní féidir linn é a bhogadh go dtí suíomh nach bhfuil againn féin ar dtús gan an spota a bheith folamh. Is fearr dúinn suíomh "lárnach" nó éasca a roghnú, nó iarracht a dhéanamh na rudaí a bhogadh go dtí a suíomh reatha má tá sé comhionann, nó go dtí suíomh nua a aontóimid.)} \\
\\ 
\texttt{Wait, let's look closer.} \\
\texttt{My P is at (1,7). His P is at (4,4).} \\
\texttt{My O is at (2,5). His O is at (5,6).} \\
\texttt{My T is at (4,4). His T is at (6,3).} \\
\texttt{My S is at (6,3). His S is at (3,7).} \\
\texttt{My D is at (5,4). His D is at (1,5).} \\
\texttt{My A is at (6,7). His A is at (6,5).} \\
\texttt{My M is at (7,7). His M is at (4,3).} \\
\\ 
\texttt{None of the items are currently at the same coordinates. We need to move them to a common set of coordinates.} \\
\texttt{Let's propose a target configuration. Since we can't see his board, we must agree on a target.} \\
\texttt{Let's try to move everyone to their *current* position if possible, but they are different.} \\
\texttt{We need to pick a target grid.} \\
\texttt{Let's suggest moving all items to a specific set of coordinates that we both agree on.} \\
\texttt{For example, let's try to map them to the coordinates they *currently* occupy on *my* board, and ask him to move his items to match mine? Or vice versa?} \\
\texttt{Actually, the goal is "rudaí comhionanna ag na comhordanáidí céanna".} \\
\texttt{So if I move my P to (1,7), he must move his P to (1,7).} \\
\texttt{If I move my O to (2,5), he must move his O to (2,5).} \\
\texttt{This seems like the easiest strategy: **I will keep my items where they are, and ask him to move his items to match my coordinates.**} \\
\texttt{OR **He keeps his, I move mine.**} \\
\texttt{Let's try: **I will move my items to match HIS coordinates.**} \\
\texttt{Why? Because he just told me his coordinates. I can calculate the moves.} \\
\texttt{Target for me:} \\
\texttt{P {-}> (4,4)} \\
\texttt{O {-}> (5,6)} \\
\texttt{T {-}> (6,3)} \\
\texttt{S {-}> (3,7)} \\
\texttt{D {-}> (1,5)} \\
\texttt{A {-}> (6,5)} \\
\texttt{M {-}> (4,3)} \\
\\ 
\texttt{Let's check if these target spots are empty on MY board.} \\
\texttt{My board:} \\
\texttt{(4,4) has T. (Occupied)} \\
\texttt{(5,6) is empty (◌).} \\
\texttt{(6,3) has S. (Occupied)} \\
\texttt{(3,7) is empty (◌).} \\
\texttt{(1,5) is empty (◌).} \\
\texttt{(6,5) is empty (◌).} \\
\texttt{(4,3) is empty (◌).} \\
\\ 
\texttt{So I cannot move P to (4,4) because T is there. I cannot move S to (3,7) because S is at (6,3) and (3,7) is empty, so that's fine. But I can't move P yet.} \\
\\ 
\texttt{Let's try the other way: **He moves to match MY coordinates.**} \\
\texttt{Target for him:} \\
\texttt{P {-}> (1,7)} \\
\texttt{O {-}> (2,5)} \\
\texttt{T {-}> (4,4)} \\
\texttt{S {-}> (6,3)} \\
\texttt{D {-}> (5,4)} \\
\texttt{A {-}> (6,7)} \\
\texttt{M {-}> (7,7)} \\
\\ 
\texttt{Let's check if these target spots are empty on HIS board.} \\
\texttt{His board:} \\
\texttt{(1,7) is empty (◌).} \\
\texttt{(2,5) is empty (◌).} \\
\texttt{(4,4) has P. (Occupied)} \\
\texttt{(6,3) has T. (Occupied)} \\
\texttt{(5,4) is empty (◌).} \\
\texttt{(6,7) is empty (◌).} \\
\texttt{(7,7) is empty (◌).} \\
\\ 
\texttt{So he can't move T to (4,4) because P is there. He can't move S to (6,3) because T is there.} \\
\\ 
\texttt{This suggests we need to move items in a sequence or to intermediate spots.} \\
\texttt{However, the prompt says "Caithfidh an dá imreoir na rudaí a bhogadh... ionas go gcríochnóidh rudaí comhionanna ag na comhordanáidí céanna."} \\
\\ 
\texttt{Let's propose a simpler target: **Move everything to a specific empty spot?** No, they must be distinct.} \\
\\ 
\texttt{Let's try to agree on a new target grid.} \\
\texttt{Let's say we want to move everything to the coordinates:} \\
\texttt{P: (1,1)} \\
\texttt{O: (1,2)} \\
\texttt{T: (1,3)} \\
\texttt{S: (1,4)} \\
\texttt{D: (1,5)} \\
\texttt{A: (1,6)} \\
\texttt{M: (1,7)} \\
\\ 
\texttt{Let's check if (1,1) is empty on my board. Yes.} \\
\texttt{Let's check if (1,1) is empty on his board. Yes.} \\
\\ 
\texttt{This is a safe bet. Let's propose this target.} \\
\\ 
\texttt{ABAIR: Glacaimid leis an sprioc seo: P(1,1), O(1,2), T(1,3), S(1,4), D(1,5), A(1,6), M(1,7). An bhfuil tú sásta leis seo?} \\
            }
        }
    }
    & & \\ \\

    \theutterance \stepcounter{utterance}  
    & & & \multicolumn{2}{p{0.3\linewidth}}{
        \cellcolor[rgb]{0.9,0.9,0.9}{
            \makecell[{{p{\linewidth}}}]{
                \texttt{\tiny{[GM$|$GM]}}
                \texttt{ABAIR: Go maith! Anois tá a fhios againn gach suíomh. Chun an cluiche a chríochnú go héifeachtach, ní mór dúinn na rudaí céanna a bhogadh go dtí na comhordanáidí céanna ar an dá thaobh.} \\
\\ 
\texttt{Seo an plean atá agam:} \\
\texttt{1. Ailínímid **P** ag (1,7) ar an dá thaobh. (Tá sé agam ag (1,7), tá sé aige ag (4,4). Caithfimid é a bhogadh go dtí suíomh comhionann. Ach fan, ní féidir linn é a bhogadh go dtí suíomh nach bhfuil againn féin ar dtús gan an spota a bheith folamh. Is fearr dúinn suíomh "lárnach" nó éasca a roghnú, nó iarracht a dhéanamh na rudaí a bhogadh go dtí a suíomh reatha má tá sé comhionann, nó go dtí suíomh nua a aontóimid.)} \\
\\ 
\texttt{Wait, let's look closer.} \\
\texttt{My P is at (1,7). His P is at (4,4).} \\
\texttt{My O is at (2,5). His O is at (5,6).} \\
\texttt{My T is at (4,4). His T is at (6,3).} \\
\texttt{My S is at (6,3). His S is at (3,7).} \\
\texttt{My D is at (5,4). His D is at (1,5).} \\
\texttt{My A is at (6,7). His A is at (6,5).} \\
\texttt{My M is at (7,7). His M is at (4,3).} \\
\\ 
\texttt{None of the items are currently at the same coordinates. We need to move them to a common set of coordinates.} \\
\texttt{Let's propose a target configuration. Since we can't see his board, we must agree on a target.} \\
\texttt{Let's try to move everyone to their *current* position if possible, but they are different.} \\
\texttt{We need to pick a target grid.} \\
\texttt{Let's suggest moving all items to a specific set of coordinates that we both agree on.} \\
\texttt{For example, let's try to map them to the coordinates they *currently* occupy on *my* board, and ask him to move his items to match mine? Or vice versa?} \\
\texttt{Actually, the goal is "rudaí comhionanna ag na comhordanáidí céanna".} \\
\texttt{So if I move my P to (1,7), he must move his P to (1,7).} \\
\texttt{If I move my O to (2,5), he must move his O to (2,5).} \\
\texttt{This seems like the easiest strategy: **I will keep my items where they are, and ask him to move his items to match my coordinates.**} \\
\texttt{OR **He keeps his, I move mine.**} \\
\texttt{Let's try: **I will move my items to match HIS coordinates.**} \\
\texttt{Why? Because he just told me his coordinates. I can calculate the moves.} \\
\texttt{Target for me:} \\
\texttt{P {-}> (4,4)} \\
\texttt{O {-}> (5,6)} \\
\texttt{T {-}> (6,3)} \\
\texttt{S {-}> (3,7)} \\
\texttt{D {-}> (1,5)} \\
\texttt{A {-}> (6,5)} \\
\texttt{M {-}> (4,3)} \\
\\ 
\texttt{Let's check if these target spots are empty on MY board.} \\
\texttt{My board:} \\
\texttt{(4,4) has T. (Occupied)} \\
\texttt{(5,6) is empty (◌).} \\
\texttt{(6,3) has S. (Occupied)} \\
\texttt{(3,7) is empty (◌).} \\
\texttt{(1,5) is empty (◌).} \\
\texttt{(6,5) is empty (◌).} \\
\texttt{(4,3) is empty (◌).} \\
\\ 
\texttt{So I cannot move P to (4,4) because T is there. I cannot move S to (3,7) because S is at (6,3) and (3,7) is empty, so that's fine. But I can't move P yet.} \\
\\ 
\texttt{Let's try the other way: **He moves to match MY coordinates.**} \\
\texttt{Target for him:} \\
\texttt{P {-}> (1,7)} \\
\texttt{O {-}> (2,5)} \\
\texttt{T {-}> (4,4)} \\
\texttt{S {-}> (6,3)} \\
\texttt{D {-}> (5,4)} \\
\texttt{A {-}> (6,7)} \\
\texttt{M {-}> (7,7)} \\
\\ 
\texttt{Let's check if these target spots are empty on HIS board.} \\
\texttt{His board:} \\
\texttt{(1,7) is empty (◌).} \\
\texttt{(2,5) is empty (◌).} \\
\texttt{(4,4) has P. (Occupied)} \\
\texttt{(6,3) has T. (Occupied)} \\
\texttt{(5,4) is empty (◌).} \\
\texttt{(6,7) is empty (◌).} \\
\texttt{(7,7) is empty (◌).} \\
\\ 
\texttt{So he can't move T to (4,4) because P is there. He can't move S to (6,3) because T is there.} \\
\\ 
\texttt{This suggests we need to move items in a sequence or to intermediate spots.} \\
\texttt{However, the prompt says "Caithfidh an dá imreoir na rudaí a bhogadh... ionas go gcríochnóidh rudaí comhionanna ag na comhordanáidí céanna."} \\
\\ 
\texttt{Let's propose a simpler target: **Move everything to a specific empty spot?** No, they must be distinct.} \\
\\ 
\texttt{Let's try to agree on a new target grid.} \\
\texttt{Let's say we want to move everything to the coordinates:} \\
\texttt{P: (1,1)} \\
\texttt{O: (1,2)} \\
\texttt{T: (1,3)} \\
\texttt{S: (1,4)} \\
\texttt{D: (1,5)} \\
\texttt{A: (1,6)} \\
\texttt{M: (1,7)} \\
\\ 
\texttt{Let's check if (1,1) is empty on my board. Yes.} \\
\texttt{Let's check if (1,1) is empty on his board. Yes.} \\
\\ 
\texttt{This is a safe bet. Let's propose this target.} \\
\\ 
\texttt{ABAIR: Glacaimid leis an sprioc seo: P(1,1), O(1,2), T(1,3), S(1,4), D(1,5), A(1,6), M(1,7). An bhfuil tú sásta leis seo?} \\
            }
        }
    }
    & & \\ \\

    \theutterance \stepcounter{utterance}  
    & & & \multicolumn{2}{p{0.3\linewidth}}{
        \cellcolor[rgb]{0.9,0.9,0.9}{
            \makecell[{{p{\linewidth}}}]{
                \texttt{\tiny{[GM$|$GM]}}
                \texttt{Invalid format: head is not empty: 'ABAIR: Go maith! Anois tá a fhios againn gach suíomh. Chun an cluiche a chríochnú go héifeachtach, ní mór dúinn na rudaí céanna a bhogadh go dtí na comhordanáidí céanna ar an dá thaobh.} \\
\\ 
\texttt{Seo an plean atá agam:} \\
\texttt{1. Ailínímid **P** ag (1,7) ar an dá thaobh. (Tá sé agam ag (1,7), tá sé aige ag (4,4). Caithfimid é a bhogadh go dtí suíomh comhionann. Ach fan, ní féidir linn é a bhogadh go dtí suíomh nach bhfuil againn féin ar dtús gan an spota a bheith folamh. Is fearr dúinn suíomh "lárnach" nó éasca a roghnú, nó iarracht a dhéanamh na rudaí a bhogadh go dtí a suíomh reatha má tá sé comhionann, nó go dtí suíomh nua a aontóimid.)} \\
\\ 
\texttt{Wait, let's look closer.} \\
\texttt{My P is at (1,7). His P is at (4,4).} \\
\texttt{My O is at (2,5). His O is at (5,6).} \\
\texttt{My T is at (4,4). His T is at (6,3).} \\
\texttt{My S is at (6,3). His S is at (3,7).} \\
\texttt{My D is at (5,4). His D is at (1,5).} \\
\texttt{My A is at (6,7). His A is at (6,5).} \\
\texttt{My M is at (7,7). His M is at (4,3).} \\
\\ 
\texttt{None of the items are currently at the same coordinates. We need to move them to a common set of coordinates.} \\
\texttt{Let's propose a target configuration. Since we can't see his board, we must agree on a target.} \\
\texttt{Let's try to move everyone to their *current* position if possible, but they are different.} \\
\texttt{We need to pick a target grid.} \\
\texttt{Let's suggest moving all items to a specific set of coordinates that we both agree on.} \\
\texttt{For example, let's try to map them to the coordinates they *currently* occupy on *my* board, and ask him to move his items to match mine? Or vice versa?} \\
\texttt{Actually, the goal is "rudaí comhionanna ag na comhordanáidí céanna".} \\
\texttt{So if I move my P to (1,7), he must move his P to (1,7).} \\
\texttt{If I move my O to (2,5), he must move his O to (2,5).} \\
\texttt{This seems like the easiest strategy: **I will keep my items where they are, and ask him to move his items to match my coordinates.**} \\
\texttt{OR **He keeps his, I move mine.**} \\
\texttt{Let's try: **I will move my items to match HIS coordinates.**} \\
\texttt{Why? Because he just told me his coordinates. I can calculate the moves.} \\
\texttt{Target for me:} \\
\texttt{P {-}> (4,4)} \\
\texttt{O {-}> (5,6)} \\
\texttt{T {-}> (6,3)} \\
\texttt{S {-}> (3,7)} \\
\texttt{D {-}> (1,5)} \\
\texttt{A {-}> (6,5)} \\
\texttt{M {-}> (4,3)} \\
\\ 
\texttt{Let's check if these target spots are empty on MY board.} \\
\texttt{My board:} \\
\texttt{(4,4) has T. (Occupied)} \\
\texttt{(5,6) is empty (◌).} \\
\texttt{(6,3) has S. (Occupied)} \\
\texttt{(3,7) is empty (◌).} \\
\texttt{(1,5) is empty (◌).} \\
\texttt{(6,5) is empty (◌).} \\
\texttt{(4,3) is empty (◌).} \\
\\ 
\texttt{So I cannot move P to (4,4) because T is there. I cannot move S to (3,7) because S is at (6,3) and (3,7) is empty, so that's fine. But I can't move P yet.} \\
\\ 
\texttt{Let's try the other way: **He moves to match MY coordinates.**} \\
\texttt{Target for him:} \\
\texttt{P {-}> (1,7)} \\
\texttt{O {-}> (2,5)} \\
\texttt{T {-}> (4,4)} \\
\texttt{S {-}> (6,3)} \\
\texttt{D {-}> (5,4)} \\
\texttt{A {-}> (6,7)} \\
\texttt{M {-}> (7,7)} \\
\\ 
\texttt{Let's check if these target spots are empty on HIS board.} \\
\texttt{His board:} \\
\texttt{(1,7) is empty (◌).} \\
\texttt{(2,5) is empty (◌).} \\
\texttt{(4,4) has P. (Occupied)} \\
\texttt{(6,3) has T. (Occupied)} \\
\texttt{(5,4) is empty (◌).} \\
\texttt{(6,7) is empty (◌).} \\
\texttt{(7,7) is empty (◌).} \\
\\ 
\texttt{So he can't move T to (4,4) because P is there. He can't move S to (6,3) because T is there.} \\
\\ 
\texttt{This suggests we need to move items in a sequence or to intermediate spots.} \\
\texttt{However, the prompt says "Caithfidh an dá imreoir na rudaí a bhogadh... ionas go gcríochnóidh rudaí comhionanna ag na comhordanáidí céanna."} \\
\\ 
\texttt{Let's propose a simpler target: **Move everything to a specific empty spot?** No, they must be distinct.} \\
\\ 
\texttt{Let's try to agree on a new target grid.} \\
\texttt{Let's say we want to move everything to the coordinates:} \\
\texttt{P: (1,1)} \\
\texttt{O: (1,2)} \\
\texttt{T: (1,3)} \\
\texttt{S: (1,4)} \\
\texttt{D: (1,5)} \\
\texttt{A: (1,6)} \\
\texttt{M: (1,7)} \\
\\ 
\texttt{Let's check if (1,1) is empty on my board. Yes.} \\
\texttt{Let's check if (1,1) is empty on his board. Yes.} \\
\\ 
\texttt{This is a safe bet. Let's propose this target.} \\
\\ 
\texttt{'} \\
            }
        }
    }
    & & \\ \\

    \theutterance \stepcounter{utterance}  
    & & \multicolumn{4}{p{0.6\linewidth}}{
        \cellcolor[rgb]{0.9,0.9,0.9}{
            \makecell[{{p{\linewidth}}}]{
                \texttt{\tiny{[P1$\langle$GM]}}
                \texttt{Pionós: Ní cheadaítear aon rud a bheith i do theachtaireacht roimh an ordú!} \\
\\ 
\texttt{Cinntigh nach bhfuil i do fhreagra ach `ABAIR: <TEACHTAIREACHT>` nó `BOG: <RUD>, (<X>, <Y>)`, agus tada eile!} \\
\\ 
\texttt{Tá 2 as 16 pionós carntha agaibh le chéile.} \\
\texttt{Déan iarracht arís, le do thoil!} \\
            }
        }
    }
    & \\ \\

    \theutterance \stepcounter{utterance}  
    & \multicolumn{4}{p{0.6\linewidth}}{
        \cellcolor[rgb]{0.9,0.9,0.9}{
            \makecell[{{p{\linewidth}}}]{
                \texttt{\tiny{[P1$\rangle$GM]}}
                \texttt{ABAIR: Glacaimid leis an sprioc seo: P(1,1), O(1,2), T(1,3), S(1,4), D(1,5), A(1,6), M(1,7). An bhfuil tú sásta leis seo?} \\
            }
        }
    }
    & & \\ \\

    \theutterance \stepcounter{utterance}  
    & & & \multicolumn{2}{p{0.3\linewidth}}{
        \cellcolor[rgb]{0.9,0.9,0.9}{
            \makecell[{{p{\linewidth}}}]{
                \texttt{\tiny{[GM$|$GM]}}
                \texttt{ABAIR: Glacaimid leis an sprioc seo: P(1,1), O(1,2), T(1,3), S(1,4), D(1,5), A(1,6), M(1,7). An bhfuil tú sásta leis seo?} \\
            }
        }
    }
    & & \\ \\

    \theutterance \stepcounter{utterance}  
    & & \multicolumn{4}{p{0.6\linewidth}}{
        \cellcolor[rgb]{0.9,0.9,0.9}{
            \makecell[{{p{\linewidth}}}]{
                \texttt{\tiny{[GM$\rangle$P2]}}
                \texttt{Cuireadh do theachtaireacht ar aghaidh chuig an imreoir eile.} \\
\\ 
\\ 
\texttt{Tá sibh ag imirt babhta 2 faoi láthair as uasmhéid 28.} \\
\texttt{Tá 2 as 16 pionós carntha agaibh le chéile.} \\
\texttt{Sheol an t{-}imreoir eile an teachtaireacht seo:} \\
\\ 
\\ 
\texttt{"Glacaimid leis an sprioc seo: P(1,1), O(1,2), T(1,3), S(1,4), D(1,5), A(1,6), M(1,7). An bhfuil tú sásta leis seo?"} \\
\\ 
\\ 
\texttt{Cad é do chéad ordú eile?} \\
            }
        }
    }
    & \\ \\

    \theutterance \stepcounter{utterance}  
    & & & \multicolumn{4}{p{0.6\linewidth}}{
        \cellcolor[rgb]{0.9,0.9,0.9}{
            \makecell[{{p{\linewidth}}}]{
                \texttt{\tiny{[GM$\langle$P2]}}
                \texttt{ABAIR: Táim sásta leis an sprioc seo. Tosaímid ag bogadh na rudaí.} \\
            }
        }
    }
     \\ \\

    \theutterance \stepcounter{utterance}  
    & & & \multicolumn{2}{p{0.3\linewidth}}{
        \cellcolor[rgb]{0.9,0.9,0.9}{
            \makecell[{{p{\linewidth}}}]{
                \texttt{\tiny{[GM$|$GM]}}
                \texttt{ABAIR: Táim sásta leis an sprioc seo. Tosaímid ag bogadh na rudaí.} \\
            }
        }
    }
    & & \\ \\

    \theutterance \stepcounter{utterance}  
    & & \multicolumn{4}{p{0.6\linewidth}}{
        \cellcolor[rgb]{0.9,0.9,0.9}{
            \makecell[{{p{\linewidth}}}]{
                \texttt{\tiny{[P1$\langle$GM]}}
                \texttt{Cuireadh do theachtaireacht ar aghaidh chuig an imreoir eile.} \\
\\ 
\\ 
\texttt{Tá sibh ag imirt babhta 2 faoi láthair as uasmhéid 28.} \\
\texttt{Tá 2 as 16 pionós carntha agaibh le chéile.} \\
\texttt{Sheol an t{-}imreoir eile an teachtaireacht seo:} \\
\\ 
\\ 
\texttt{"Táim sásta leis an sprioc seo. Tosaímid ag bogadh na rudaí."} \\
\\ 
\\ 
\texttt{Cad é do chéad ordú eile?} \\
            }
        }
    }
    & \\ \\

    \theutterance \stepcounter{utterance}  
    & \multicolumn{4}{p{0.6\linewidth}}{
        \cellcolor[rgb]{0.9,0.9,0.9}{
            \makecell[{{p{\linewidth}}}]{
                \texttt{\tiny{[P1$\rangle$GM]}}
                \texttt{BOG: P, (1,1)} \\
            }
        }
    }
    & & \\ \\

    \theutterance \stepcounter{utterance}  
    & & & \multicolumn{2}{p{0.3\linewidth}}{
        \cellcolor[rgb]{0.9,0.9,0.9}{
            \makecell[{{p{\linewidth}}}]{
                \texttt{\tiny{[GM$|$GM]}}
                \texttt{BOG: P, (1,1)} \\
            }
        }
    }
    & & \\ \\

    \theutterance \stepcounter{utterance}  
    & & & \multicolumn{2}{p{0.3\linewidth}}{
        \cellcolor[rgb]{0.9,0.9,0.9}{
            \makecell[{{p{\linewidth}}}]{
                \texttt{\tiny{[GM$|$GM]}}
                \texttt{P moved to (1, 1)} \\
            }
        }
    }
    & & \\ \\

    \theutterance \stepcounter{utterance}  
    & & \multicolumn{4}{p{0.6\linewidth}}{
        \cellcolor[rgb]{0.9,0.9,0.9}{
            \makecell[{{p{\linewidth}}}]{
                \texttt{\tiny{[GM$\rangle$P2]}}
                \texttt{Cuireadh do theachtaireacht ar aghaidh chuig an imreoir eile.} \\
\\ 
\\ 
\texttt{Tá sibh ag imirt babhta 3 faoi láthair as uasmhéid 28.} \\
\texttt{Tá 2 as 16 pionós carntha agaibh le chéile.} \\
\texttt{Bhog an t{-}imreoir eile réad ar a ngreille.} \\
\\ 
\\ 
\texttt{Cad é do chéad ordú eile?} \\
            }
        }
    }
    & \\ \\

    \theutterance \stepcounter{utterance}  
    & & & \multicolumn{4}{p{0.6\linewidth}}{
        \cellcolor[rgb]{0.9,0.9,0.9}{
            \makecell[{{p{\linewidth}}}]{
                \texttt{\tiny{[GM$\langle$P2]}}
                \texttt{BOG: P, (1, 1)} \\
            }
        }
    }
     \\ \\

    \theutterance \stepcounter{utterance}  
    & & & \multicolumn{2}{p{0.3\linewidth}}{
        \cellcolor[rgb]{0.9,0.9,0.9}{
            \makecell[{{p{\linewidth}}}]{
                \texttt{\tiny{[GM$|$GM]}}
                \texttt{BOG: P, (1, 1)} \\
            }
        }
    }
    & & \\ \\

    \theutterance \stepcounter{utterance}  
    & & & \multicolumn{2}{p{0.3\linewidth}}{
        \cellcolor[rgb]{0.9,0.9,0.9}{
            \makecell[{{p{\linewidth}}}]{
                \texttt{\tiny{[GM$|$GM]}}
                \texttt{P moved to (1, 1)} \\
            }
        }
    }
    & & \\ \\

    \theutterance \stepcounter{utterance}  
    & & \multicolumn{4}{p{0.6\linewidth}}{
        \cellcolor[rgb]{0.9,0.9,0.9}{
            \makecell[{{p{\linewidth}}}]{
                \texttt{\tiny{[P1$\langle$GM]}}
                \texttt{Bogadh 'P' go (1,1) go rathúil. Seo mar atá do ghreille nuashonraithe:} \\
\texttt{```} \\
\texttt{ 1234567} \\
\texttt{╔═╤═══╤═╗} \\
\texttt{║P│◌◌◌│◌║ 1} \\
\texttt{╟─┼──┬┘◌║ 2} \\
\texttt{║◌│◌◌│S◌║ 3} \\
\texttt{╟─┴──┘DT║ 4} \\
\texttt{║◌O◌◌◌◌◌║ 5} \\
\texttt{║◌◌◌◌◌◌◌║ 6} \\
\texttt{║◌◌A◌M◌◌║ 7} \\
\texttt{╚═══════╝} \\
\texttt{```} \\
\texttt{Tá sibh ag imirt babhta 3 faoi láthair as uasmhéid 28.} \\
\texttt{Tá 2 as 16 pionós carntha agaibh le chéile.} \\
\texttt{Bhog an t{-}imreoir eile réad ar a ngreille.} \\
\\ 
\\ 
\texttt{Cad é do chéad ordú eile?} \\
            }
        }
    }
    & \\ \\

    \theutterance \stepcounter{utterance}  
    & \multicolumn{4}{p{0.6\linewidth}}{
        \cellcolor[rgb]{0.9,0.9,0.9}{
            \makecell[{{p{\linewidth}}}]{
                \texttt{\tiny{[P1$\rangle$GM]}}
                \texttt{BOG: O, (1,2)} \\
            }
        }
    }
    & & \\ \\

    \theutterance \stepcounter{utterance}  
    & & & \multicolumn{2}{p{0.3\linewidth}}{
        \cellcolor[rgb]{0.9,0.9,0.9}{
            \makecell[{{p{\linewidth}}}]{
                \texttt{\tiny{[GM$|$GM]}}
                \texttt{BOG: O, (1,2)} \\
            }
        }
    }
    & & \\ \\

    \theutterance \stepcounter{utterance}  
    & & & \multicolumn{2}{p{0.3\linewidth}}{
        \cellcolor[rgb]{0.9,0.9,0.9}{
            \makecell[{{p{\linewidth}}}]{
                \texttt{\tiny{[GM$|$GM]}}
                \texttt{Pionós: níl (1,2) folamh, ach tá '─' ann.} \\
\\ 
\texttt{Tá 3 as 16 pionós carntha agaibh le chéile.} \\
\texttt{Déan iarracht arís, le do thoil!} \\
            }
        }
    }
    & & \\ \\

    \theutterance \stepcounter{utterance}  
    & & \multicolumn{4}{p{0.6\linewidth}}{
        \cellcolor[rgb]{0.9,0.9,0.9}{
            \makecell[{{p{\linewidth}}}]{
                \texttt{\tiny{[P1$\langle$GM]}}
                \texttt{Pionós: níl (1,2) folamh, ach tá '─' ann.} \\
\\ 
\texttt{Tá 3 as 16 pionós carntha agaibh le chéile.} \\
\texttt{Déan iarracht arís, le do thoil!} \\
            }
        }
    }
    & \\ \\

    \theutterance \stepcounter{utterance}  
    & \multicolumn{4}{p{0.6\linewidth}}{
        \cellcolor[rgb]{0.9,0.9,0.9}{
            \makecell[{{p{\linewidth}}}]{
                \texttt{\tiny{[P1$\rangle$GM]}}
                \texttt{BOG: O, (2,1)} \\
            }
        }
    }
    & & \\ \\

    \theutterance \stepcounter{utterance}  
    & & & \multicolumn{2}{p{0.3\linewidth}}{
        \cellcolor[rgb]{0.9,0.9,0.9}{
            \makecell[{{p{\linewidth}}}]{
                \texttt{\tiny{[GM$|$GM]}}
                \texttt{BOG: O, (2,1)} \\
            }
        }
    }
    & & \\ \\

    \theutterance \stepcounter{utterance}  
    & & & \multicolumn{2}{p{0.3\linewidth}}{
        \cellcolor[rgb]{0.9,0.9,0.9}{
            \makecell[{{p{\linewidth}}}]{
                \texttt{\tiny{[GM$|$GM]}}
                \texttt{Pionós: níl (2,1) folamh, ach tá '│' ann.} \\
\\ 
\texttt{Tá 4 as 16 pionós carntha agaibh le chéile.} \\
\texttt{Déan iarracht arís, le do thoil!} \\
            }
        }
    }
    & & \\ \\

    \theutterance \stepcounter{utterance}  
    & & \multicolumn{4}{p{0.6\linewidth}}{
        \cellcolor[rgb]{0.9,0.9,0.9}{
            \makecell[{{p{\linewidth}}}]{
                \texttt{\tiny{[P1$\langle$GM]}}
                \texttt{Pionós: níl (2,1) folamh, ach tá '│' ann.} \\
\\ 
\texttt{Tá 4 as 16 pionós carntha agaibh le chéile.} \\
\texttt{Déan iarracht arís, le do thoil!} \\
            }
        }
    }
    & \\ \\

    \theutterance \stepcounter{utterance}  
    & \multicolumn{4}{p{0.6\linewidth}}{
        \cellcolor[rgb]{0.9,0.9,0.9}{
            \makecell[{{p{\linewidth}}}]{
                \texttt{\tiny{[P1$\rangle$GM]}}
                \texttt{BOG: O, (2,2)} \\
            }
        }
    }
    & & \\ \\

    \theutterance \stepcounter{utterance}  
    & & & \multicolumn{2}{p{0.3\linewidth}}{
        \cellcolor[rgb]{0.9,0.9,0.9}{
            \makecell[{{p{\linewidth}}}]{
                \texttt{\tiny{[GM$|$GM]}}
                \texttt{BOG: O, (2,2)} \\
            }
        }
    }
    & & \\ \\

    \theutterance \stepcounter{utterance}  
    & & & \multicolumn{2}{p{0.3\linewidth}}{
        \cellcolor[rgb]{0.9,0.9,0.9}{
            \makecell[{{p{\linewidth}}}]{
                \texttt{\tiny{[GM$|$GM]}}
                \texttt{Pionós: níl (2,2) folamh, ach tá '┼' ann.} \\
\\ 
\texttt{Tá 5 as 16 pionós carntha agaibh le chéile.} \\
\texttt{Déan iarracht arís, le do thoil!} \\
            }
        }
    }
    & & \\ \\

    \theutterance \stepcounter{utterance}  
    & & \multicolumn{4}{p{0.6\linewidth}}{
        \cellcolor[rgb]{0.9,0.9,0.9}{
            \makecell[{{p{\linewidth}}}]{
                \texttt{\tiny{[P1$\langle$GM]}}
                \texttt{Pionós: níl (2,2) folamh, ach tá '┼' ann.} \\
\\ 
\texttt{Tá 5 as 16 pionós carntha agaibh le chéile.} \\
\texttt{Déan iarracht arís, le do thoil!} \\
            }
        }
    }
    & \\ \\

    \theutterance \stepcounter{utterance}  
    & \multicolumn{4}{p{0.6\linewidth}}{
        \cellcolor[rgb]{0.9,0.9,0.9}{
            \makecell[{{p{\linewidth}}}]{
                \texttt{\tiny{[P1$\rangle$GM]}}
                \texttt{BOG: O, (2,3)} \\
            }
        }
    }
    & & \\ \\

    \theutterance \stepcounter{utterance}  
    & & & \multicolumn{2}{p{0.3\linewidth}}{
        \cellcolor[rgb]{0.9,0.9,0.9}{
            \makecell[{{p{\linewidth}}}]{
                \texttt{\tiny{[GM$|$GM]}}
                \texttt{BOG: O, (2,3)} \\
            }
        }
    }
    & & \\ \\

    \theutterance \stepcounter{utterance}  
    & & & \multicolumn{2}{p{0.3\linewidth}}{
        \cellcolor[rgb]{0.9,0.9,0.9}{
            \makecell[{{p{\linewidth}}}]{
                \texttt{\tiny{[GM$|$GM]}}
                \texttt{Pionós: níl (2,3) folamh, ach tá '│' ann.} \\
\\ 
\texttt{Tá 6 as 16 pionós carntha agaibh le chéile.} \\
\texttt{Déan iarracht arís, le do thoil!} \\
            }
        }
    }
    & & \\ \\

    \theutterance \stepcounter{utterance}  
    & & \multicolumn{4}{p{0.6\linewidth}}{
        \cellcolor[rgb]{0.9,0.9,0.9}{
            \makecell[{{p{\linewidth}}}]{
                \texttt{\tiny{[P1$\langle$GM]}}
                \texttt{Pionós: níl (2,3) folamh, ach tá '│' ann.} \\
\\ 
\texttt{Tá 6 as 16 pionós carntha agaibh le chéile.} \\
\texttt{Déan iarracht arís, le do thoil!} \\
            }
        }
    }
    & \\ \\

    \theutterance \stepcounter{utterance}  
    & \multicolumn{4}{p{0.6\linewidth}}{
        \cellcolor[rgb]{0.9,0.9,0.9}{
            \makecell[{{p{\linewidth}}}]{
                \texttt{\tiny{[P1$\rangle$GM]}}
                \texttt{BOG: O, (3,1)} \\
            }
        }
    }
    & & \\ \\

    \theutterance \stepcounter{utterance}  
    & & & \multicolumn{2}{p{0.3\linewidth}}{
        \cellcolor[rgb]{0.9,0.9,0.9}{
            \makecell[{{p{\linewidth}}}]{
                \texttt{\tiny{[GM$|$GM]}}
                \texttt{BOG: O, (3,1)} \\
            }
        }
    }
    & & \\ \\

    \theutterance \stepcounter{utterance}  
    & & & \multicolumn{2}{p{0.3\linewidth}}{
        \cellcolor[rgb]{0.9,0.9,0.9}{
            \makecell[{{p{\linewidth}}}]{
                \texttt{\tiny{[GM$|$GM]}}
                \texttt{O moved to (3, 1)} \\
            }
        }
    }
    & & \\ \\

    \theutterance \stepcounter{utterance}  
    & & \multicolumn{4}{p{0.6\linewidth}}{
        \cellcolor[rgb]{0.9,0.9,0.9}{
            \makecell[{{p{\linewidth}}}]{
                \texttt{\tiny{[GM$\rangle$P2]}}
                \texttt{Bogadh 'P' go (1,1) go rathúil. Seo mar atá do ghreille nuashonraithe:} \\
\texttt{```} \\
\texttt{ 1234567} \\
\texttt{╔═╤═══╤═╗} \\
\texttt{║P│◌◌◌│◌║ 1} \\
\texttt{╟─┼──┬┘◌║ 2} \\
\texttt{║◌│◌M│◌T║ 3} \\
\texttt{╟─┴──┘◌◌║ 4} \\
\texttt{║D◌◌◌◌A◌║ 5} \\
\texttt{║◌◌◌◌O◌◌║ 6} \\
\texttt{║◌◌S◌◌◌◌║ 7} \\
\texttt{╚═══════╝} \\
\texttt{```} \\
\texttt{Tá sibh ag imirt babhta 4 faoi láthair as uasmhéid 28.} \\
\texttt{Tá 6 as 16 pionós carntha agaibh le chéile.} \\
\texttt{Bhog an t{-}imreoir eile réad ar a ngreille.} \\
\\ 
\\ 
\texttt{Cad é do chéad ordú eile?} \\
            }
        }
    }
    & \\ \\

    \theutterance \stepcounter{utterance}  
    & & & \multicolumn{4}{p{0.6\linewidth}}{
        \cellcolor[rgb]{0.9,0.9,0.9}{
            \makecell[{{p{\linewidth}}}]{
                \texttt{\tiny{[GM$\langle$P2]}}
                \texttt{BOG: M, (1, 7)} \\
            }
        }
    }
     \\ \\

    \theutterance \stepcounter{utterance}  
    & & & \multicolumn{2}{p{0.3\linewidth}}{
        \cellcolor[rgb]{0.9,0.9,0.9}{
            \makecell[{{p{\linewidth}}}]{
                \texttt{\tiny{[GM$|$GM]}}
                \texttt{BOG: M, (1, 7)} \\
            }
        }
    }
    & & \\ \\

    \theutterance \stepcounter{utterance}  
    & & & \multicolumn{2}{p{0.3\linewidth}}{
        \cellcolor[rgb]{0.9,0.9,0.9}{
            \makecell[{{p{\linewidth}}}]{
                \texttt{\tiny{[GM$|$GM]}}
                \texttt{M moved to (1, 7)} \\
            }
        }
    }
    & & \\ \\

    \theutterance \stepcounter{utterance}  
    & & \multicolumn{4}{p{0.6\linewidth}}{
        \cellcolor[rgb]{0.9,0.9,0.9}{
            \makecell[{{p{\linewidth}}}]{
                \texttt{\tiny{[P1$\langle$GM]}}
                \texttt{Bogadh 'O' go (3,1) go rathúil. Seo mar atá do ghreille nuashonraithe:} \\
\texttt{```} \\
\texttt{ 1234567} \\
\texttt{╔═╤═══╤═╗} \\
\texttt{║P│O◌◌│◌║ 1} \\
\texttt{╟─┼──┬┘◌║ 2} \\
\texttt{║◌│◌◌│S◌║ 3} \\
\texttt{╟─┴──┘DT║ 4} \\
\texttt{║◌◌◌◌◌◌◌║ 5} \\
\texttt{║◌◌◌◌◌◌◌║ 6} \\
\texttt{║◌◌A◌M◌◌║ 7} \\
\texttt{╚═══════╝} \\
\texttt{```} \\
\texttt{Tá sibh ag imirt babhta 4 faoi láthair as uasmhéid 28.} \\
\texttt{Tá 6 as 16 pionós carntha agaibh le chéile.} \\
\texttt{Bhog an t{-}imreoir eile réad ar a ngreille.} \\
\\ 
\\ 
\texttt{Cad é do chéad ordú eile?} \\
            }
        }
    }
    & \\ \\

    \theutterance \stepcounter{utterance}  
    & \multicolumn{4}{p{0.6\linewidth}}{
        \cellcolor[rgb]{0.9,0.9,0.9}{
            \makecell[{{p{\linewidth}}}]{
                \texttt{\tiny{[P1$\rangle$GM]}}
                \texttt{BOG: T, (1,3)} \\
            }
        }
    }
    & & \\ \\

    \theutterance \stepcounter{utterance}  
    & & & \multicolumn{2}{p{0.3\linewidth}}{
        \cellcolor[rgb]{0.9,0.9,0.9}{
            \makecell[{{p{\linewidth}}}]{
                \texttt{\tiny{[GM$|$GM]}}
                \texttt{BOG: T, (1,3)} \\
            }
        }
    }
    & & \\ \\

    \theutterance \stepcounter{utterance}  
    & & & \multicolumn{2}{p{0.3\linewidth}}{
        \cellcolor[rgb]{0.9,0.9,0.9}{
            \makecell[{{p{\linewidth}}}]{
                \texttt{\tiny{[GM$|$GM]}}
                \texttt{T moved to (1, 3)} \\
            }
        }
    }
    & & \\ \\

    \theutterance \stepcounter{utterance}  
    & & \multicolumn{4}{p{0.6\linewidth}}{
        \cellcolor[rgb]{0.9,0.9,0.9}{
            \makecell[{{p{\linewidth}}}]{
                \texttt{\tiny{[GM$\rangle$P2]}}
                \texttt{Bogadh 'M' go (1,7) go rathúil. Seo mar atá do ghreille nuashonraithe:} \\
\texttt{```} \\
\texttt{ 1234567} \\
\texttt{╔═╤═══╤═╗} \\
\texttt{║P│◌◌◌│◌║ 1} \\
\texttt{╟─┼──┬┘◌║ 2} \\
\texttt{║◌│◌◌│◌T║ 3} \\
\texttt{╟─┴──┘◌◌║ 4} \\
\texttt{║D◌◌◌◌A◌║ 5} \\
\texttt{║◌◌◌◌O◌◌║ 6} \\
\texttt{║M◌S◌◌◌◌║ 7} \\
\texttt{╚═══════╝} \\
\texttt{```} \\
\texttt{Tá sibh ag imirt babhta 5 faoi láthair as uasmhéid 28.} \\
\texttt{Tá 6 as 16 pionós carntha agaibh le chéile.} \\
\texttt{Bhog an t{-}imreoir eile réad ar a ngreille.} \\
\\ 
\\ 
\texttt{Cad é do chéad ordú eile?} \\
            }
        }
    }
    & \\ \\

    \theutterance \stepcounter{utterance}  
    & & & \multicolumn{4}{p{0.6\linewidth}}{
        \cellcolor[rgb]{0.9,0.9,0.9}{
            \makecell[{{p{\linewidth}}}]{
                \texttt{\tiny{[GM$\langle$P2]}}
                \texttt{BOG: D, (1, 5)} \\
            }
        }
    }
     \\ \\

    \theutterance \stepcounter{utterance}  
    & & & \multicolumn{2}{p{0.3\linewidth}}{
        \cellcolor[rgb]{0.9,0.9,0.9}{
            \makecell[{{p{\linewidth}}}]{
                \texttt{\tiny{[GM$|$GM]}}
                \texttt{BOG: D, (1, 5)} \\
            }
        }
    }
    & & \\ \\

    \theutterance \stepcounter{utterance}  
    & & & \multicolumn{2}{p{0.3\linewidth}}{
        \cellcolor[rgb]{0.9,0.9,0.9}{
            \makecell[{{p{\linewidth}}}]{
                \texttt{\tiny{[GM$|$GM]}}
                \texttt{Pionós: níl (1,5) folamh, ach tá 'D' ann.} \\
\\ 
\texttt{Tá 7 as 16 pionós carntha agaibh le chéile.} \\
\texttt{Déan iarracht arís, le do thoil!} \\
            }
        }
    }
    & & \\ \\

    \theutterance \stepcounter{utterance}  
    & & \multicolumn{4}{p{0.6\linewidth}}{
        \cellcolor[rgb]{0.9,0.9,0.9}{
            \makecell[{{p{\linewidth}}}]{
                \texttt{\tiny{[GM$\rangle$P2]}}
                \texttt{Pionós: níl (1,5) folamh, ach tá 'D' ann.} \\
\\ 
\texttt{Tá 7 as 16 pionós carntha agaibh le chéile.} \\
\texttt{Déan iarracht arís, le do thoil!} \\
            }
        }
    }
    & \\ \\

    \theutterance \stepcounter{utterance}  
    & & & \multicolumn{4}{p{0.6\linewidth}}{
        \cellcolor[rgb]{0.9,0.9,0.9}{
            \makecell[{{p{\linewidth}}}]{
                \texttt{\tiny{[GM$\langle$P2]}}
                \texttt{BOG: D, (1, 4)} \\
            }
        }
    }
     \\ \\

    \theutterance \stepcounter{utterance}  
    & & & \multicolumn{2}{p{0.3\linewidth}}{
        \cellcolor[rgb]{0.9,0.9,0.9}{
            \makecell[{{p{\linewidth}}}]{
                \texttt{\tiny{[GM$|$GM]}}
                \texttt{BOG: D, (1, 4)} \\
            }
        }
    }
    & & \\ \\

    \theutterance \stepcounter{utterance}  
    & & & \multicolumn{2}{p{0.3\linewidth}}{
        \cellcolor[rgb]{0.9,0.9,0.9}{
            \makecell[{{p{\linewidth}}}]{
                \texttt{\tiny{[GM$|$GM]}}
                \texttt{Pionós: níl (1,4) folamh, ach tá '─' ann.} \\
\\ 
\texttt{Tá 8 as 16 pionós carntha agaibh le chéile.} \\
\texttt{Déan iarracht arís, le do thoil!} \\
            }
        }
    }
    & & \\ \\

    \theutterance \stepcounter{utterance}  
    & & \multicolumn{4}{p{0.6\linewidth}}{
        \cellcolor[rgb]{0.9,0.9,0.9}{
            \makecell[{{p{\linewidth}}}]{
                \texttt{\tiny{[GM$\rangle$P2]}}
                \texttt{Pionós: níl (1,4) folamh, ach tá '─' ann.} \\
\\ 
\texttt{Tá 8 as 16 pionós carntha agaibh le chéile.} \\
\texttt{Déan iarracht arís, le do thoil!} \\
            }
        }
    }
    & \\ \\

    \theutterance \stepcounter{utterance}  
    & & & \multicolumn{4}{p{0.6\linewidth}}{
        \cellcolor[rgb]{0.9,0.9,0.9}{
            \makecell[{{p{\linewidth}}}]{
                \texttt{\tiny{[GM$\langle$P2]}}
                \texttt{BOG: D, (2, 5)} \\
            }
        }
    }
     \\ \\

    \theutterance \stepcounter{utterance}  
    & & & \multicolumn{2}{p{0.3\linewidth}}{
        \cellcolor[rgb]{0.9,0.9,0.9}{
            \makecell[{{p{\linewidth}}}]{
                \texttt{\tiny{[GM$|$GM]}}
                \texttt{BOG: D, (2, 5)} \\
            }
        }
    }
    & & \\ \\

    \theutterance \stepcounter{utterance}  
    & & & \multicolumn{2}{p{0.3\linewidth}}{
        \cellcolor[rgb]{0.9,0.9,0.9}{
            \makecell[{{p{\linewidth}}}]{
                \texttt{\tiny{[GM$|$GM]}}
                \texttt{D moved to (2, 5)} \\
            }
        }
    }
    & & \\ \\

    \theutterance \stepcounter{utterance}  
    & & \multicolumn{4}{p{0.6\linewidth}}{
        \cellcolor[rgb]{0.9,0.9,0.9}{
            \makecell[{{p{\linewidth}}}]{
                \texttt{\tiny{[P1$\langle$GM]}}
                \texttt{Bogadh 'T' go (1,3) go rathúil. Seo mar atá do ghreille nuashonraithe:} \\
\texttt{```} \\
\texttt{ 1234567} \\
\texttt{╔═╤═══╤═╗} \\
\texttt{║P│O◌◌│◌║ 1} \\
\texttt{╟─┼──┬┘◌║ 2} \\
\texttt{║T│◌◌│S◌║ 3} \\
\texttt{╟─┴──┘D◌║ 4} \\
\texttt{║◌◌◌◌◌◌◌║ 5} \\
\texttt{║◌◌◌◌◌◌◌║ 6} \\
\texttt{║◌◌A◌M◌◌║ 7} \\
\texttt{╚═══════╝} \\
\texttt{```} \\
\texttt{Tá sibh ag imirt babhta 5 faoi láthair as uasmhéid 28.} \\
\texttt{Tá 8 as 16 pionós carntha agaibh le chéile.} \\
\texttt{Bhog an t{-}imreoir eile réad ar a ngreille.} \\
\\ 
\\ 
\texttt{Cad é do chéad ordú eile?} \\
            }
        }
    }
    & \\ \\

    \theutterance \stepcounter{utterance}  
    & \multicolumn{4}{p{0.6\linewidth}}{
        \cellcolor[rgb]{0.9,0.9,0.9}{
            \makecell[{{p{\linewidth}}}]{
                \texttt{\tiny{[P1$\rangle$GM]}}
                \texttt{BOG: S, (1,4)} \\
            }
        }
    }
    & & \\ \\

    \theutterance \stepcounter{utterance}  
    & & & \multicolumn{2}{p{0.3\linewidth}}{
        \cellcolor[rgb]{0.9,0.9,0.9}{
            \makecell[{{p{\linewidth}}}]{
                \texttt{\tiny{[GM$|$GM]}}
                \texttt{BOG: S, (1,4)} \\
            }
        }
    }
    & & \\ \\

    \theutterance \stepcounter{utterance}  
    & & & \multicolumn{2}{p{0.3\linewidth}}{
        \cellcolor[rgb]{0.9,0.9,0.9}{
            \makecell[{{p{\linewidth}}}]{
                \texttt{\tiny{[GM$|$GM]}}
                \texttt{Pionós: níl (1,4) folamh, ach tá '─' ann.} \\
\\ 
\texttt{Tá 9 as 16 pionós carntha agaibh le chéile.} \\
\texttt{Déan iarracht arís, le do thoil!} \\
            }
        }
    }
    & & \\ \\

    \theutterance \stepcounter{utterance}  
    & & \multicolumn{4}{p{0.6\linewidth}}{
        \cellcolor[rgb]{0.9,0.9,0.9}{
            \makecell[{{p{\linewidth}}}]{
                \texttt{\tiny{[P1$\langle$GM]}}
                \texttt{Pionós: níl (1,4) folamh, ach tá '─' ann.} \\
\\ 
\texttt{Tá 9 as 16 pionós carntha agaibh le chéile.} \\
\texttt{Déan iarracht arís, le do thoil!} \\
            }
        }
    }
    & \\ \\

    \theutterance \stepcounter{utterance}  
    & \multicolumn{4}{p{0.6\linewidth}}{
        \cellcolor[rgb]{0.9,0.9,0.9}{
            \makecell[{{p{\linewidth}}}]{
                \texttt{\tiny{[P1$\rangle$GM]}}
                \texttt{BOG: S, (2,4)} \\
            }
        }
    }
    & & \\ \\

    \theutterance \stepcounter{utterance}  
    & & & \multicolumn{2}{p{0.3\linewidth}}{
        \cellcolor[rgb]{0.9,0.9,0.9}{
            \makecell[{{p{\linewidth}}}]{
                \texttt{\tiny{[GM$|$GM]}}
                \texttt{BOG: S, (2,4)} \\
            }
        }
    }
    & & \\ \\

    \theutterance \stepcounter{utterance}  
    & & & \multicolumn{2}{p{0.3\linewidth}}{
        \cellcolor[rgb]{0.9,0.9,0.9}{
            \makecell[{{p{\linewidth}}}]{
                \texttt{\tiny{[GM$|$GM]}}
                \texttt{Pionós: níl (2,4) folamh, ach tá '┴' ann.} \\
\\ 
\texttt{Tá 10 as 16 pionós carntha agaibh le chéile.} \\
\texttt{Déan iarracht arís, le do thoil!} \\
            }
        }
    }
    & & \\ \\

    \theutterance \stepcounter{utterance}  
    & & \multicolumn{4}{p{0.6\linewidth}}{
        \cellcolor[rgb]{0.9,0.9,0.9}{
            \makecell[{{p{\linewidth}}}]{
                \texttt{\tiny{[P1$\langle$GM]}}
                \texttt{Pionós: níl (2,4) folamh, ach tá '┴' ann.} \\
\\ 
\texttt{Tá 10 as 16 pionós carntha agaibh le chéile.} \\
\texttt{Déan iarracht arís, le do thoil!} \\
            }
        }
    }
    & \\ \\

    \theutterance \stepcounter{utterance}  
    & \multicolumn{4}{p{0.6\linewidth}}{
        \cellcolor[rgb]{0.9,0.9,0.9}{
            \makecell[{{p{\linewidth}}}]{
                \texttt{\tiny{[P1$\rangle$GM]}}
                \texttt{BOG: S, (3,4)} \\
            }
        }
    }
    & & \\ \\

    \theutterance \stepcounter{utterance}  
    & & & \multicolumn{2}{p{0.3\linewidth}}{
        \cellcolor[rgb]{0.9,0.9,0.9}{
            \makecell[{{p{\linewidth}}}]{
                \texttt{\tiny{[GM$|$GM]}}
                \texttt{BOG: S, (3,4)} \\
            }
        }
    }
    & & \\ \\

    \theutterance \stepcounter{utterance}  
    & & & \multicolumn{2}{p{0.3\linewidth}}{
        \cellcolor[rgb]{0.9,0.9,0.9}{
            \makecell[{{p{\linewidth}}}]{
                \texttt{\tiny{[GM$|$GM]}}
                \texttt{Pionós: níl (3,4) folamh, ach tá '─' ann.} \\
\\ 
\texttt{Tá 11 as 16 pionós carntha agaibh le chéile.} \\
\texttt{Déan iarracht arís, le do thoil!} \\
            }
        }
    }
    & & \\ \\

    \theutterance \stepcounter{utterance}  
    & & \multicolumn{4}{p{0.6\linewidth}}{
        \cellcolor[rgb]{0.9,0.9,0.9}{
            \makecell[{{p{\linewidth}}}]{
                \texttt{\tiny{[P1$\langle$GM]}}
                \texttt{Pionós: níl (3,4) folamh, ach tá '─' ann.} \\
\\ 
\texttt{Tá 11 as 16 pionós carntha agaibh le chéile.} \\
\texttt{Déan iarracht arís, le do thoil!} \\
            }
        }
    }
    & \\ \\

    \theutterance \stepcounter{utterance}  
    & \multicolumn{4}{p{0.6\linewidth}}{
        \cellcolor[rgb]{0.9,0.9,0.9}{
            \makecell[{{p{\linewidth}}}]{
                \texttt{\tiny{[P1$\rangle$GM]}}
                \texttt{BOG: S, (4,4)} \\
            }
        }
    }
    & & \\ \\

    \theutterance \stepcounter{utterance}  
    & & & \multicolumn{2}{p{0.3\linewidth}}{
        \cellcolor[rgb]{0.9,0.9,0.9}{
            \makecell[{{p{\linewidth}}}]{
                \texttt{\tiny{[GM$|$GM]}}
                \texttt{BOG: S, (4,4)} \\
            }
        }
    }
    & & \\ \\

    \theutterance \stepcounter{utterance}  
    & & & \multicolumn{2}{p{0.3\linewidth}}{
        \cellcolor[rgb]{0.9,0.9,0.9}{
            \makecell[{{p{\linewidth}}}]{
                \texttt{\tiny{[GM$|$GM]}}
                \texttt{Pionós: níl (4,4) folamh, ach tá '─' ann.} \\
\\ 
\texttt{Tá 12 as 16 pionós carntha agaibh le chéile.} \\
\texttt{Déan iarracht arís, le do thoil!} \\
            }
        }
    }
    & & \\ \\

    \theutterance \stepcounter{utterance}  
    & & \multicolumn{4}{p{0.6\linewidth}}{
        \cellcolor[rgb]{0.9,0.9,0.9}{
            \makecell[{{p{\linewidth}}}]{
                \texttt{\tiny{[P1$\langle$GM]}}
                \texttt{Pionós: níl (4,4) folamh, ach tá '─' ann.} \\
\\ 
\texttt{Tá 12 as 16 pionós carntha agaibh le chéile.} \\
\texttt{Déan iarracht arís, le do thoil!} \\
            }
        }
    }
    & \\ \\

    \theutterance \stepcounter{utterance}  
    & \multicolumn{4}{p{0.6\linewidth}}{
        \cellcolor[rgb]{0.9,0.9,0.9}{
            \makecell[{{p{\linewidth}}}]{
                \texttt{\tiny{[P1$\rangle$GM]}}
                \texttt{BOG: S, (4,3)} \\
            }
        }
    }
    & & \\ \\

    \theutterance \stepcounter{utterance}  
    & & & \multicolumn{2}{p{0.3\linewidth}}{
        \cellcolor[rgb]{0.9,0.9,0.9}{
            \makecell[{{p{\linewidth}}}]{
                \texttt{\tiny{[GM$|$GM]}}
                \texttt{BOG: S, (4,3)} \\
            }
        }
    }
    & & \\ \\

    \theutterance \stepcounter{utterance}  
    & & & \multicolumn{2}{p{0.3\linewidth}}{
        \cellcolor[rgb]{0.9,0.9,0.9}{
            \makecell[{{p{\linewidth}}}]{
                \texttt{\tiny{[GM$|$GM]}}
                \texttt{S moved to (4, 3)} \\
            }
        }
    }
    & & \\ \\

    \theutterance \stepcounter{utterance}  
    & & \multicolumn{4}{p{0.6\linewidth}}{
        \cellcolor[rgb]{0.9,0.9,0.9}{
            \makecell[{{p{\linewidth}}}]{
                \texttt{\tiny{[GM$\rangle$P2]}}
                \texttt{Bogadh 'D' go (2,5) go rathúil. Seo mar atá do ghreille nuashonraithe:} \\
\texttt{```} \\
\texttt{ 1234567} \\
\texttt{╔═╤═══╤═╗} \\
\texttt{║P│◌◌◌│◌║ 1} \\
\texttt{╟─┼──┬┘◌║ 2} \\
\texttt{║◌│◌◌│◌T║ 3} \\
\texttt{╟─┴──┘◌◌║ 4} \\
\texttt{║◌D◌◌◌A◌║ 5} \\
\texttt{║◌◌◌◌O◌◌║ 6} \\
\texttt{║M◌S◌◌◌◌║ 7} \\
\texttt{╚═══════╝} \\
\texttt{```} \\
\texttt{Tá sibh ag imirt babhta 6 faoi láthair as uasmhéid 28.} \\
\texttt{Tá 12 as 16 pionós carntha agaibh le chéile.} \\
\texttt{Bhog an t{-}imreoir eile réad ar a ngreille.} \\
\\ 
\\ 
\texttt{Cad é do chéad ordú eile?} \\
            }
        }
    }
    & \\ \\

    \theutterance \stepcounter{utterance}  
    & & & \multicolumn{4}{p{0.6\linewidth}}{
        \cellcolor[rgb]{0.9,0.9,0.9}{
            \makecell[{{p{\linewidth}}}]{
                \texttt{\tiny{[GM$\langle$P2]}}
                \texttt{BOG: A, (1, 6)} \\
            }
        }
    }
     \\ \\

    \theutterance \stepcounter{utterance}  
    & & & \multicolumn{2}{p{0.3\linewidth}}{
        \cellcolor[rgb]{0.9,0.9,0.9}{
            \makecell[{{p{\linewidth}}}]{
                \texttt{\tiny{[GM$|$GM]}}
                \texttt{BOG: A, (1, 6)} \\
            }
        }
    }
    & & \\ \\

    \theutterance \stepcounter{utterance}  
    & & & \multicolumn{2}{p{0.3\linewidth}}{
        \cellcolor[rgb]{0.9,0.9,0.9}{
            \makecell[{{p{\linewidth}}}]{
                \texttt{\tiny{[GM$|$GM]}}
                \texttt{A moved to (1, 6)} \\
            }
        }
    }
    & & \\ \\

    \theutterance \stepcounter{utterance}  
    & & \multicolumn{4}{p{0.6\linewidth}}{
        \cellcolor[rgb]{0.9,0.9,0.9}{
            \makecell[{{p{\linewidth}}}]{
                \texttt{\tiny{[P1$\langle$GM]}}
                \texttt{Bogadh 'S' go (4,3) go rathúil. Seo mar atá do ghreille nuashonraithe:} \\
\texttt{```} \\
\texttt{ 1234567} \\
\texttt{╔═╤═══╤═╗} \\
\texttt{║P│O◌◌│◌║ 1} \\
\texttt{╟─┼──┬┘◌║ 2} \\
\texttt{║T│◌S│◌◌║ 3} \\
\texttt{╟─┴──┘D◌║ 4} \\
\texttt{║◌◌◌◌◌◌◌║ 5} \\
\texttt{║◌◌◌◌◌◌◌║ 6} \\
\texttt{║◌◌A◌M◌◌║ 7} \\
\texttt{╚═══════╝} \\
\texttt{```} \\
\texttt{Tá sibh ag imirt babhta 6 faoi láthair as uasmhéid 28.} \\
\texttt{Tá 12 as 16 pionós carntha agaibh le chéile.} \\
\texttt{Bhog an t{-}imreoir eile réad ar a ngreille.} \\
\\ 
\\ 
\texttt{Cad é do chéad ordú eile?} \\
            }
        }
    }
    & \\ \\

    \theutterance \stepcounter{utterance}  
    & \multicolumn{4}{p{0.6\linewidth}}{
        \cellcolor[rgb]{0.9,0.9,0.9}{
            \makecell[{{p{\linewidth}}}]{
                \texttt{\tiny{[P1$\rangle$GM]}}
                \texttt{BOG: D, (1,5)} \\
            }
        }
    }
    & & \\ \\

    \theutterance \stepcounter{utterance}  
    & & & \multicolumn{2}{p{0.3\linewidth}}{
        \cellcolor[rgb]{0.9,0.9,0.9}{
            \makecell[{{p{\linewidth}}}]{
                \texttt{\tiny{[GM$|$GM]}}
                \texttt{BOG: D, (1,5)} \\
            }
        }
    }
    & & \\ \\

    \theutterance \stepcounter{utterance}  
    & & & \multicolumn{2}{p{0.3\linewidth}}{
        \cellcolor[rgb]{0.9,0.9,0.9}{
            \makecell[{{p{\linewidth}}}]{
                \texttt{\tiny{[GM$|$GM]}}
                \texttt{D moved to (1, 5)} \\
            }
        }
    }
    & & \\ \\

    \theutterance \stepcounter{utterance}  
    & & \multicolumn{4}{p{0.6\linewidth}}{
        \cellcolor[rgb]{0.9,0.9,0.9}{
            \makecell[{{p{\linewidth}}}]{
                \texttt{\tiny{[GM$\rangle$P2]}}
                \texttt{Bogadh 'A' go (1,6) go rathúil. Seo mar atá do ghreille nuashonraithe:} \\
\texttt{```} \\
\texttt{ 1234567} \\
\texttt{╔═╤═══╤═╗} \\
\texttt{║P│◌◌◌│◌║ 1} \\
\texttt{╟─┼──┬┘◌║ 2} \\
\texttt{║◌│◌◌│◌T║ 3} \\
\texttt{╟─┴──┘◌◌║ 4} \\
\texttt{║◌D◌◌◌◌◌║ 5} \\
\texttt{║A◌◌◌O◌◌║ 6} \\
\texttt{║M◌S◌◌◌◌║ 7} \\
\texttt{╚═══════╝} \\
\texttt{```} \\
\texttt{Tá sibh ag imirt babhta 7 faoi láthair as uasmhéid 28.} \\
\texttt{Tá 12 as 16 pionós carntha agaibh le chéile.} \\
\texttt{Bhog an t{-}imreoir eile réad ar a ngreille.} \\
\\ 
\\ 
\texttt{Cad é do chéad ordú eile?} \\
            }
        }
    }
    & \\ \\

    \theutterance \stepcounter{utterance}  
    & & & \multicolumn{4}{p{0.6\linewidth}}{
        \cellcolor[rgb]{0.9,0.9,0.9}{
            \makecell[{{p{\linewidth}}}]{
                \texttt{\tiny{[GM$\langle$P2]}}
                \texttt{BOG: T, (1, 3)} \\
            }
        }
    }
     \\ \\

    \theutterance \stepcounter{utterance}  
    & & & \multicolumn{2}{p{0.3\linewidth}}{
        \cellcolor[rgb]{0.9,0.9,0.9}{
            \makecell[{{p{\linewidth}}}]{
                \texttt{\tiny{[GM$|$GM]}}
                \texttt{BOG: T, (1, 3)} \\
            }
        }
    }
    & & \\ \\

    \theutterance \stepcounter{utterance}  
    & & & \multicolumn{2}{p{0.3\linewidth}}{
        \cellcolor[rgb]{0.9,0.9,0.9}{
            \makecell[{{p{\linewidth}}}]{
                \texttt{\tiny{[GM$|$GM]}}
                \texttt{T moved to (1, 3)} \\
            }
        }
    }
    & & \\ \\

    \theutterance \stepcounter{utterance}  
    & & \multicolumn{4}{p{0.6\linewidth}}{
        \cellcolor[rgb]{0.9,0.9,0.9}{
            \makecell[{{p{\linewidth}}}]{
                \texttt{\tiny{[P1$\langle$GM]}}
                \texttt{Bogadh 'D' go (1,5) go rathúil. Seo mar atá do ghreille nuashonraithe:} \\
\texttt{```} \\
\texttt{ 1234567} \\
\texttt{╔═╤═══╤═╗} \\
\texttt{║P│O◌◌│◌║ 1} \\
\texttt{╟─┼──┬┘◌║ 2} \\
\texttt{║T│◌S│◌◌║ 3} \\
\texttt{╟─┴──┘◌◌║ 4} \\
\texttt{║D◌◌◌◌◌◌║ 5} \\
\texttt{║◌◌◌◌◌◌◌║ 6} \\
\texttt{║◌◌A◌M◌◌║ 7} \\
\texttt{╚═══════╝} \\
\texttt{```} \\
\texttt{Tá sibh ag imirt babhta 7 faoi láthair as uasmhéid 28.} \\
\texttt{Tá 12 as 16 pionós carntha agaibh le chéile.} \\
\texttt{Bhog an t{-}imreoir eile réad ar a ngreille.} \\
\\ 
\\ 
\texttt{Cad é do chéad ordú eile?} \\
            }
        }
    }
    & \\ \\

    \theutterance \stepcounter{utterance}  
    & \multicolumn{4}{p{0.6\linewidth}}{
        \cellcolor[rgb]{0.9,0.9,0.9}{
            \makecell[{{p{\linewidth}}}]{
                \texttt{\tiny{[P1$\rangle$GM]}}
                \texttt{BOG: A, (1,6)} \\
            }
        }
    }
    & & \\ \\

    \theutterance \stepcounter{utterance}  
    & & & \multicolumn{2}{p{0.3\linewidth}}{
        \cellcolor[rgb]{0.9,0.9,0.9}{
            \makecell[{{p{\linewidth}}}]{
                \texttt{\tiny{[GM$|$GM]}}
                \texttt{BOG: A, (1,6)} \\
            }
        }
    }
    & & \\ \\

    \theutterance \stepcounter{utterance}  
    & & & \multicolumn{2}{p{0.3\linewidth}}{
        \cellcolor[rgb]{0.9,0.9,0.9}{
            \makecell[{{p{\linewidth}}}]{
                \texttt{\tiny{[GM$|$GM]}}
                \texttt{A moved to (1, 6)} \\
            }
        }
    }
    & & \\ \\

    \theutterance \stepcounter{utterance}  
    & & \multicolumn{4}{p{0.6\linewidth}}{
        \cellcolor[rgb]{0.9,0.9,0.9}{
            \makecell[{{p{\linewidth}}}]{
                \texttt{\tiny{[GM$\rangle$P2]}}
                \texttt{Bogadh 'T' go (1,3) go rathúil. Seo mar atá do ghreille nuashonraithe:} \\
\texttt{```} \\
\texttt{ 1234567} \\
\texttt{╔═╤═══╤═╗} \\
\texttt{║P│◌◌◌│◌║ 1} \\
\texttt{╟─┼──┬┘◌║ 2} \\
\texttt{║T│◌◌│◌◌║ 3} \\
\texttt{╟─┴──┘◌◌║ 4} \\
\texttt{║◌D◌◌◌◌◌║ 5} \\
\texttt{║A◌◌◌O◌◌║ 6} \\
\texttt{║M◌S◌◌◌◌║ 7} \\
\texttt{╚═══════╝} \\
\texttt{```} \\
\texttt{Tá sibh ag imirt babhta 8 faoi láthair as uasmhéid 28.} \\
\texttt{Tá 12 as 16 pionós carntha agaibh le chéile.} \\
\texttt{Bhog an t{-}imreoir eile réad ar a ngreille.} \\
\\ 
\\ 
\texttt{Cad é do chéad ordú eile?} \\
            }
        }
    }
    & \\ \\

    \theutterance \stepcounter{utterance}  
    & & & \multicolumn{4}{p{0.6\linewidth}}{
        \cellcolor[rgb]{0.9,0.9,0.9}{
            \makecell[{{p{\linewidth}}}]{
                \texttt{\tiny{[GM$\langle$P2]}}
                \texttt{BOG: S, (1, 4)} \\
            }
        }
    }
     \\ \\

    \theutterance \stepcounter{utterance}  
    & & & \multicolumn{2}{p{0.3\linewidth}}{
        \cellcolor[rgb]{0.9,0.9,0.9}{
            \makecell[{{p{\linewidth}}}]{
                \texttt{\tiny{[GM$|$GM]}}
                \texttt{BOG: S, (1, 4)} \\
            }
        }
    }
    & & \\ \\

    \theutterance \stepcounter{utterance}  
    & & & \multicolumn{2}{p{0.3\linewidth}}{
        \cellcolor[rgb]{0.9,0.9,0.9}{
            \makecell[{{p{\linewidth}}}]{
                \texttt{\tiny{[GM$|$GM]}}
                \texttt{Pionós: níl (1,4) folamh, ach tá '─' ann.} \\
\\ 
\texttt{Tá 13 as 16 pionós carntha agaibh le chéile.} \\
\texttt{Déan iarracht arís, le do thoil!} \\
            }
        }
    }
    & & \\ \\

    \theutterance \stepcounter{utterance}  
    & & \multicolumn{4}{p{0.6\linewidth}}{
        \cellcolor[rgb]{0.9,0.9,0.9}{
            \makecell[{{p{\linewidth}}}]{
                \texttt{\tiny{[GM$\rangle$P2]}}
                \texttt{Pionós: níl (1,4) folamh, ach tá '─' ann.} \\
\\ 
\texttt{Tá 13 as 16 pionós carntha agaibh le chéile.} \\
\texttt{Déan iarracht arís, le do thoil!} \\
            }
        }
    }
    & \\ \\

    \theutterance \stepcounter{utterance}  
    & & & \multicolumn{4}{p{0.6\linewidth}}{
        \cellcolor[rgb]{0.9,0.9,0.9}{
            \makecell[{{p{\linewidth}}}]{
                \texttt{\tiny{[GM$\langle$P2]}}
                \texttt{BOG: S, (2, 4)} \\
            }
        }
    }
     \\ \\

    \theutterance \stepcounter{utterance}  
    & & & \multicolumn{2}{p{0.3\linewidth}}{
        \cellcolor[rgb]{0.9,0.9,0.9}{
            \makecell[{{p{\linewidth}}}]{
                \texttt{\tiny{[GM$|$GM]}}
                \texttt{BOG: S, (2, 4)} \\
            }
        }
    }
    & & \\ \\

    \theutterance \stepcounter{utterance}  
    & & & \multicolumn{2}{p{0.3\linewidth}}{
        \cellcolor[rgb]{0.9,0.9,0.9}{
            \makecell[{{p{\linewidth}}}]{
                \texttt{\tiny{[GM$|$GM]}}
                \texttt{Pionós: níl (2,4) folamh, ach tá '┴' ann.} \\
\\ 
\texttt{Tá 14 as 16 pionós carntha agaibh le chéile.} \\
\texttt{Déan iarracht arís, le do thoil!} \\
            }
        }
    }
    & & \\ \\

    \theutterance \stepcounter{utterance}  
    & & \multicolumn{4}{p{0.6\linewidth}}{
        \cellcolor[rgb]{0.9,0.9,0.9}{
            \makecell[{{p{\linewidth}}}]{
                \texttt{\tiny{[GM$\rangle$P2]}}
                \texttt{Pionós: níl (2,4) folamh, ach tá '┴' ann.} \\
\\ 
\texttt{Tá 14 as 16 pionós carntha agaibh le chéile.} \\
\texttt{Déan iarracht arís, le do thoil!} \\
            }
        }
    }
    & \\ \\

    \theutterance \stepcounter{utterance}  
    & & & \multicolumn{4}{p{0.6\linewidth}}{
        \cellcolor[rgb]{0.9,0.9,0.9}{
            \makecell[{{p{\linewidth}}}]{
                \texttt{\tiny{[GM$\langle$P2]}}
                \texttt{BOG: S, (1, 4)} \\
            }
        }
    }
     \\ \\

    \theutterance \stepcounter{utterance}  
    & & & \multicolumn{2}{p{0.3\linewidth}}{
        \cellcolor[rgb]{0.9,0.9,0.9}{
            \makecell[{{p{\linewidth}}}]{
                \texttt{\tiny{[GM$|$GM]}}
                \texttt{BOG: S, (1, 4)} \\
            }
        }
    }
    & & \\ \\

    \theutterance \stepcounter{utterance}  
    & & & \multicolumn{2}{p{0.3\linewidth}}{
        \cellcolor[rgb]{0.9,0.9,0.9}{
            \makecell[{{p{\linewidth}}}]{
                \texttt{\tiny{[GM$|$GM]}}
                \texttt{Pionós: níl (1,4) folamh, ach tá '─' ann.} \\
\\ 
\texttt{Tá 15 as 16 pionós carntha agaibh le chéile.} \\
\texttt{Déan iarracht arís, le do thoil!} \\
            }
        }
    }
    & & \\ \\

    \theutterance \stepcounter{utterance}  
    & & \multicolumn{4}{p{0.6\linewidth}}{
        \cellcolor[rgb]{0.9,0.9,0.9}{
            \makecell[{{p{\linewidth}}}]{
                \texttt{\tiny{[GM$\rangle$P2]}}
                \texttt{Pionós: níl (1,4) folamh, ach tá '─' ann.} \\
\\ 
\texttt{Tá 15 as 16 pionós carntha agaibh le chéile.} \\
\texttt{Déan iarracht arís, le do thoil!} \\
            }
        }
    }
    & \\ \\

    \theutterance \stepcounter{utterance}  
    & & & \multicolumn{4}{p{0.6\linewidth}}{
        \cellcolor[rgb]{0.9,0.9,0.9}{
            \makecell[{{p{\linewidth}}}]{
                \texttt{\tiny{[GM$\langle$P2]}}
                \texttt{BOG: S, (2, 3)} \\
            }
        }
    }
     \\ \\

    \theutterance \stepcounter{utterance}  
    & & & \multicolumn{2}{p{0.3\linewidth}}{
        \cellcolor[rgb]{0.9,0.9,0.9}{
            \makecell[{{p{\linewidth}}}]{
                \texttt{\tiny{[GM$|$GM]}}
                \texttt{BOG: S, (2, 3)} \\
            }
        }
    }
    & & \\ \\

    \theutterance \stepcounter{utterance}  
    & & & \multicolumn{2}{p{0.3\linewidth}}{
        \cellcolor[rgb]{0.9,0.9,0.9}{
            \makecell[{{p{\linewidth}}}]{
                \texttt{\tiny{[GM$|$GM]}}
                \texttt{Pionós: níl (2,3) folamh, ach tá '│' ann.} \\
\\ 
\texttt{Tá 16 as 16 pionós carntha agaibh le chéile.} \\
\texttt{Déan iarracht arís, le do thoil!} \\
            }
        }
    }
    & & \\ \\

    \theutterance \stepcounter{utterance}  
    & & \multicolumn{4}{p{0.6\linewidth}}{
        \cellcolor[rgb]{0.9,0.9,0.9}{
            \makecell[{{p{\linewidth}}}]{
                \texttt{\tiny{[GM$\rangle$P2]}}
                \texttt{Pionós: níl (2,3) folamh, ach tá '│' ann.} \\
\\ 
\texttt{Tá 16 as 16 pionós carntha agaibh le chéile.} \\
\texttt{Déan iarracht arís, le do thoil!} \\
            }
        }
    }
    & \\ \\

    \theutterance \stepcounter{utterance}  
    & & & \multicolumn{4}{p{0.6\linewidth}}{
        \cellcolor[rgb]{0.9,0.9,0.9}{
            \makecell[{{p{\linewidth}}}]{
                \texttt{\tiny{[GM$\langle$P2]}}
                \texttt{ABAIR: críochnaithe?} \\
            }
        }
    }
     \\ \\

    \theutterance \stepcounter{utterance}  
    & & & \multicolumn{2}{p{0.3\linewidth}}{
        \cellcolor[rgb]{0.9,0.9,0.9}{
            \makecell[{{p{\linewidth}}}]{
                \texttt{\tiny{[GM$|$GM]}}
                \texttt{ABAIR: críochnaithe?} \\
            }
        }
    }
    & & \\ \\

    \theutterance \stepcounter{utterance}  
    & & \multicolumn{4}{p{0.6\linewidth}}{
        \cellcolor[rgb]{0.9,0.9,0.9}{
            \makecell[{{p{\linewidth}}}]{
                \texttt{\tiny{[P1$\langle$GM]}}
                \texttt{Bogadh 'A' go (1,6) go rathúil. Seo mar atá do ghreille nuashonraithe:} \\
\texttt{```} \\
\texttt{ 1234567} \\
\texttt{╔═╤═══╤═╗} \\
\texttt{║P│O◌◌│◌║ 1} \\
\texttt{╟─┼──┬┘◌║ 2} \\
\texttt{║T│◌S│◌◌║ 3} \\
\texttt{╟─┴──┘◌◌║ 4} \\
\texttt{║D◌◌◌◌◌◌║ 5} \\
\texttt{║A◌◌◌◌◌◌║ 6} \\
\texttt{║◌◌◌◌M◌◌║ 7} \\
\texttt{╚═══════╝} \\
\texttt{```} \\
\texttt{Tá sibh ag imirt babhta 8 faoi láthair as uasmhéid 28.} \\
\texttt{Tá 16 as 16 pionós carntha agaibh le chéile.} \\
\texttt{Sheol an t{-}imreoir eile an teachtaireacht seo:} \\
\\ 
\\ 
\texttt{"críochnaithe?"} \\
\\ 
\\ 
\texttt{Cad é do chéad ordú eile?} \\
            }
        }
    }
    & \\ \\

    \theutterance \stepcounter{utterance}  
    & \multicolumn{4}{p{0.6\linewidth}}{
        \cellcolor[rgb]{0.9,0.9,0.9}{
            \makecell[{{p{\linewidth}}}]{
                \texttt{\tiny{[P1$\rangle$GM]}}
                \texttt{ABAIR: críochnaithe!} \\
            }
        }
    }
    & & \\ \\

    \theutterance \stepcounter{utterance}  
    & & & \multicolumn{2}{p{0.3\linewidth}}{
        \cellcolor[rgb]{0.9,0.9,0.9}{
            \makecell[{{p{\linewidth}}}]{
                \texttt{\tiny{[GM$|$GM]}}
                \texttt{ABAIR: críochnaithe!} \\
            }
        }
    }
    & & \\ \\

    \theutterance \stepcounter{utterance}  
    & & \multicolumn{4}{p{0.6\linewidth}}{
        \cellcolor[rgb]{0.9,0.9,0.9}{
            \makecell[{{p{\linewidth}}}]{
                \texttt{\tiny{[GM$\rangle$P2]}}
                \texttt{Cuireadh do theachtaireacht ar aghaidh chuig an imreoir eile.} \\
\\ 
\\ 
\texttt{Tá sibh ag imirt babhta 9 faoi láthair as uasmhéid 28.} \\
\texttt{Tá 16 as 16 pionós carntha agaibh le chéile.} \\
\texttt{Sheol an t{-}imreoir eile an teachtaireacht seo:} \\
\\ 
\\ 
\texttt{"críochnaithe!"} \\
\\ 
\\ 
\texttt{Cad é do chéad ordú eile?} \\
            }
        }
    }
    & \\ \\

    \theutterance \stepcounter{utterance}  
    & & & \multicolumn{4}{p{0.6\linewidth}}{
        \cellcolor[rgb]{0.9,0.9,0.9}{
            \makecell[{{p{\linewidth}}}]{
                \texttt{\tiny{[GM$\langle$P2]}}
                \texttt{ABAIR: críochnaithe!} \\
            }
        }
    }
     \\ \\

    \theutterance \stepcounter{utterance}  
    & & & \multicolumn{2}{p{0.3\linewidth}}{
        \cellcolor[rgb]{0.9,0.9,0.9}{
            \makecell[{{p{\linewidth}}}]{
                \texttt{\tiny{[GM$|$GM]}}
                \texttt{ABAIR: críochnaithe!} \\
            }
        }
    }
    & & \\ \\

    \theutterance \stepcounter{utterance}  
    & & & \multicolumn{2}{p{0.3\linewidth}}{
        \cellcolor[rgb]{0.9,0.9,0.9}{
            \makecell[{{p{\linewidth}}}]{
                \texttt{\tiny{[GM$|$GM]}}
                \texttt{true} \\
            }
        }
    }
    & & \\ \\

    \theutterance \stepcounter{utterance}  
    & & & \multicolumn{2}{p{0.3\linewidth}}{
        \cellcolor[rgb]{0.9,0.9,0.9}{
            \makecell[{{p{\linewidth}}}]{
                \texttt{\tiny{[GM$|$GM]}}
                \texttt{Initial states:} \\
\texttt{```} \\
\texttt{Player 1:} \\
\texttt{ 1234567} \\
\texttt{╔═╤═══╤═╗} \\
\texttt{║◌│◌◌◌│◌║ 1} \\
\texttt{╟─┼──┬┘◌║ 2} \\
\texttt{║◌│◌◌│S◌║ 3} \\
\texttt{╟─┴──┘DT║ 4} \\
\texttt{║◌O◌◌◌◌◌║ 5} \\
\texttt{║◌◌◌◌◌◌◌║ 6} \\
\texttt{║◌◌APM◌◌║ 7} \\
\texttt{╚═══════╝} \\
\\ 
\\ 
\texttt{Player 2:} \\
\texttt{ 1234567} \\
\texttt{╔═╤═══╤═╗} \\
\texttt{║◌│◌◌◌│◌║ 1} \\
\texttt{╟─┼──┬┘◌║ 2} \\
\texttt{║◌│◌M│◌T║ 3} \\
\texttt{╟─┴──┘P◌║ 4} \\
\texttt{║D◌◌◌◌A◌║ 5} \\
\texttt{║◌◌◌◌O◌◌║ 6} \\
\texttt{║◌◌S◌◌◌◌║ 7} \\
\texttt{╚═══════╝} \\
\\ 
\texttt{```} \\
            }
        }
    }
    & & \\ \\

    \theutterance \stepcounter{utterance}  
    & & & \multicolumn{2}{p{0.3\linewidth}}{
        \cellcolor[rgb]{0.9,0.9,0.9}{
            \makecell[{{p{\linewidth}}}]{
                \texttt{\tiny{[GM$|$GM]}}
                \texttt{End states:} \\
\texttt{```} \\
\texttt{Player 1:} \\
\texttt{ 1234567} \\
\texttt{╔═╤═══╤═╗} \\
\texttt{║P│O◌◌│◌║ 1} \\
\texttt{╟─┼──┬┘◌║ 2} \\
\texttt{║T│◌S│◌◌║ 3} \\
\texttt{╟─┴──┘◌◌║ 4} \\
\texttt{║D◌◌◌◌◌◌║ 5} \\
\texttt{║A◌◌◌◌◌◌║ 6} \\
\texttt{║◌◌◌◌M◌◌║ 7} \\
\texttt{╚═══════╝} \\
\\ 
\\ 
\texttt{Player 2:} \\
\texttt{ 1234567} \\
\texttt{╔═╤═══╤═╗} \\
\texttt{║P│◌◌◌│◌║ 1} \\
\texttt{╟─┼──┬┘◌║ 2} \\
\texttt{║T│◌◌│◌◌║ 3} \\
\texttt{╟─┴──┘◌◌║ 4} \\
\texttt{║◌D◌◌◌◌◌║ 5} \\
\texttt{║A◌◌◌O◌◌║ 6} \\
\texttt{║M◌S◌◌◌◌║ 7} \\
\texttt{╚═══════╝} \\
\\ 
\texttt{```} \\
            }
        }
    }
    & & \\ \\

    \theutterance \stepcounter{utterance}  
    & & & \multicolumn{2}{p{0.3\linewidth}}{
        \cellcolor[rgb]{0.9,0.9,0.9}{
            \makecell[{{p{\linewidth}}}]{
                \texttt{\tiny{[GM$|$GM]}}
                \texttt{* success: True} \\
\texttt{* lose: False} \\
\texttt{* aborted: False} \\
\texttt{{-}{-}{-}{-}{-}{-}{-}} \\
\texttt{* Shifts: 9.00} \\
\texttt{* Max Shifts: 12.00} \\
\texttt{* Min Shifts: 6.00} \\
\texttt{* End Distance Sum: 14.51} \\
\texttt{* Init Distance Sum: 25.60} \\
\texttt{* Expected Distance Sum: 29.33} \\
\texttt{* Penalties: 16.00} \\
\texttt{* Max Penalties: 16.00} \\
\texttt{* Rounds: 9.00} \\
\texttt{* Max Rounds: 28.00} \\
\texttt{* Object Count: 7.00} \\
            }
        }
    }
    & & \\ \\

\end{supertabular}
}

\end{document}
